% SPDX-FileCopyrightText: 2026 Will Cook
% SPDX-License-Identifier: CC-BY-4.0
%
% Long-form reasoning record seeded from the complete problem note.
% The short note and this record are now independent publication units.
% This flat manuscript is generated from reviewable parts below.
\documentclass[11pt]{article}
% SPDX-FileCopyrightText: 2026 Will Cook
% SPDX-License-Identifier: CC-BY-4.0
%
% Shared preamble for the Erdős Problem Notes: the short, standalone,
% problem-owned manuscripts covering the expansion library `ErdosProblems`.
%
% One note owns one Erdős problem.  Each note repeats whatever context its own
% reader needs, so that a reader who arrives at a single problem never has to
% assemble it from the gateway paper.  Deliberate repetition across notes is the
% design, not an oversight.
%
% Source links resolve against one pinned formal-source commit, declared once
% below.  `scripts/check_problem_note_sources.py` verifies that every authored
% (file, line, declaration) triple names that exact declaration in the pinned
% snapshot, so a link cannot silently rot as later work moves lines.

\usepackage{paper-house-style}
\usepackage{array}
\usepackage{booktabs}
\usepackage{enumitem}
\setlist{topsep=0.35em,itemsep=0.2em}
\usepackage[colorlinks=true,linkcolor=linkinternal,urlcolor=linkexternal,
  citecolor=linkinternal]{hyperref}
\hypersetup{bookmarksnumbered=true,bookmarksopen=true,bookmarksopenlevel=1,
  pdfauthor={Will Cook},pdflang={en-GB}}

% ---------------------------------------------------------------------------
% Pinned formal source.  \commit names the checkpoint every link resolves
% against; the checker reads that snapshot rather than the working tree, so the
% printed coordinates stay correct however later waves move declarations.
% ---------------------------------------------------------------------------
\newcommand{\commit}{e5fd26a6d1b1a346430205eab5385b4c9565d004}
\newcommand{\commitshort}{e5fd26a6d1b1}
\newcommand{\repobase}{https://github.com/wcook04/plectis-lean-erdos249-257/blob/\commit}
\newcommand{\sourceurl}{https://github.com/wcook04/plectis-lean-erdos249-257/tree/\commit}
\newcommand{\PX}{ErdosProblems}

% Declaration names stay inside link targets.  The printed label is either a
% spaced, lower-case reading of the declaration name (\lref) or a short authored
% phrase (\lword); neither prints a module path or a raw Lean identifier.
\ExplSyntaxOn
\cs_new_protected:Npn \noteleanlabel #1
  {
    \tl_set:Nx \l_tmpa_tl { \tl_to_str:n {#1} }
    \regex_replace_all:nnN { _+ } { \c{space} } \l_tmpa_tl
    \regex_replace_all:nnN
      { ([a-z0-9])([A-Z]) }
      { \1\c{space}\2 }
      \l_tmpa_tl
    \tl_set:Nx \l_tmpa_tl { \tl_lower_case:n { \tl_use:N \l_tmpa_tl } }
    \tl_use:N \l_tmpa_tl
  }
\ExplSyntaxOff
\newcommand{\leanlabel}[1]{\noteleanlabel{#1}}
\newcommand{\lref}[3]{\href{\repobase/\PX/#1\#L#2}{\leanlabel{#3}}}
\newcommand{\lword}[4]{\href{\repobase/\PX/#1\#L#2}{#4}}
\newcommand{\lloc}[2]{\href{\repobase/\PX/#1\#L#2}{Lean source}}

% Links into a sibling library, named repository-relative.  The expansion
% library is implicit in \lref/\lword, because that is where problem-owned
% work lands.  Several problems have their headline theorem in the older
% reviewed #249/#257 corpus, and a note that could not name that library
% would have to paraphrase a checked statement instead of linking it.
\newcommand{\mref}[3]{\href{\repobase/#1\#L#2}{\leanlabel{#3}}}
\newcommand{\mword}[4]{\href{\repobase/#1\#L#2}{#4}}
\newcommand{\mloc}[2]{\href{\repobase/#1\#L#2}{Lean source}}
\emergencystretch=2em

\newtheorem{theorem}{Theorem}[section]
\newtheorem{proposition}[theorem]{Proposition}
\newtheorem{lemma}[theorem]{Lemma}
\newtheorem{corollary}[theorem]{Corollary}
\theoremstyle{definition}
\newtheorem{definition}[theorem]{Definition}
\newtheorem{example}[theorem]{Example}
\newtheorem{problem}[theorem]{Problem}
\theoremstyle{remark}
\newtheorem*{remark}{Remark}

\newcommand{\N}{\mathbb{N}}
\newcommand{\Npos}{\mathbb{N}_{>0}}
\newcommand{\Z}{\mathbb{Z}}
\newcommand{\Q}{\mathbb{Q}}
\newcommand{\R}{\mathbb{R}}
\newcommand{\lcm}{\operatorname{lcm}}
\newcommand{\ph}{\varphi}

% The series line printed under every note's title block.
\newcommand{\seriesline}{Erd\H{o}s Problem Notes}

% Production provenance.  The wording is fixed by the corpus provenance
% contract and reproduced verbatim in every manuscript in this repository, so
% the papers cannot drift into different accounts of how they were made.
\newcommand{\noteprovenance}{\provenancenote{\emph{Authorship and AI use.} The
prose, formal proofs, and software in this version were drafted and revised by
large-language-model agents under Will Cook's direction.  Cook set the
objectives and acceptance criteria, reviewed and authorised the manuscript, and
is responsible for its contents.  The agents are production tools, not authors.
See \emph{Statements and declarations}.}}

% The status paragraph every note carries, in its own words, before any result.
% Kernel checking establishes the exact Lean proposition; it does not establish
% that the proposition is the interesting one, that it is new, or that the
% Erdős problem is closed.
\newcommand{\statusboundary}{%
  \emph{Status.}  The problem treated here is open, and this note does not
  close it.  Every statement below marked as checked is a proposition that the
  pinned Lean kernel accepts from the sources this note links to, with no
  \texttt{sorry}, no added axiom, and no unchecked evaluation.  That is a claim
  about the formal statement, not about its mathematical interest, its
  novelty, or the original problem.  The unresolved obligations are named
  exactly, in their own section, and none of the finite computations,
  reductions, or no-go results here removes one of them.}

% Compatibility labels for statements moved from the short paper.
% amsmath permits only one ordinary label per equation. Its preserved kernel
% label writer lets the aliases share the same number and PDF anchor.
\makeatletter
\newcommand{\compatlabel}[1]{\ltx@label{#1}}
\makeatother
\usepackage[htt]{hyphenat}
\usepackage{newunicodechar}
\newunicodechar{ℕ}{\ensuremath{\N}}
\newunicodechar{ℚ}{\ensuremath{\Q}}
\newunicodechar{ℤ}{\ensuremath{\Z}}
\newunicodechar{₀}{\ensuremath{{}_0}}
% Problem-specific source pin: #251 advances independently of the corpus
% default, while the #249/#257 notes retain their own separately pinned snapshot.
\renewcommand{\commit}{3d6d938d696fed0fb71dd55115a18a73738ff223}
\renewcommand{\commitshort}{3d6d938d696f}
\hypersetup{
  pdftitle={Prime-gap dyadic series: perturbations, exact criteria and certificates},
  pdfsubject={Erdős Problem Notes: Problem 251},
  pdfkeywords={irrationality, prime gaps, dyadic series, summation by parts, Lean 4},
}

\runninghead{Erd\H{o}s \#251: perturbations, criteria and certificates}
\title{Prime-Gap Dyadic Series:\\Perturbations, Exact Criteria and Certificates}
\subtitle{A complete reasoning record for Erd\H{o}s Problem~\#251}
\author{Will Cook}
\date{September 2026}

\setlength{\emergencystretch}{2em}
\def\UrlBreaks{\do\.\do\_\do\/}
\begin{document}
% ---- part core ----
\maketitle
\noteprovenance

\begin{abstract}
We study what information about an integer coefficient word can determine
irrationality of its dyadic sum.  Our main obstruction is a sparse
perturbation theorem: for any summable nonnegative integer word and any
$f(n)\to\infty$, one fixed support of upper Banach density zero realises
an interval of target values using corrections eventually bounded by $f$,
while preserving every fixed eventual congruence of coefficients and
cumulative sums.  Polylogarithmic corrections also preserve the empirical
laws of literal blocks of length $o(\log\log X)$.  These conclusions do not
preserve primality.

For the actual zero-based primes and gaps we prove
$\Pi=\sum_{n\ge0}p_n2^{-(n+1)}=2+\sum_{n\ge0}g_n2^{-(n+1)}$ and identify the
complete tails $T_N=\sum_{j\ge1}g_{N+j}2^{-j}$.  Their recurrence gives exact
integral-shift criteria for irrationality.  Joint small shifts with unequal
associated gaps would suffice, but their required cofinal occurrence is
not proved.  Explicit countermodels separate this missing prime-specific
input from growth, parity, congruence and fixed-block nonconcentration.
We give exact finite window certificates and rational-denominator
exclusions, distinguishing ordinary proofs, independently replayable
integer computations and the statements covered by the supplied formal
verification records.  No unconditional irrationality result for the
prime series is claimed.
\end{abstract}

\section{Introduction}\label{long251:sec:problem}

Throughout, the primes are indexed from zero: $p_0=2$, $p_1=3$,
$p_2=5$, and $g_n=p_{n+1}-p_n$.  Erd\H{o}s Problem~\#251 concerns
\[
 \Pi=\sum_{n\ge0}\frac{p_n}{2^{n+1}}
     =2/2+3/4+5/8+7/16+\cdots,
\]
with the question recorded in \cite[p.~94]{erdos1958},
\cite[p.~62]{erdosgraham1980} and \cite[p.~103]{erdos1988}.
The supplied catalogue snapshot records no unconditional solution
\cite{erdosproblems}; no refreshed status of the live catalogue is assumed
in this revision.  The first values are
\[
 \begin{array}{c|cccccccccc}
  i   & 0&1&2&3&4 &5 &6 &7 &8 &9\\\hline
  p_i & 2&3&5&7&11&13&17&19&23&29\\
  g_i & 1&2&2&4&2 &4 &2 &4 &6 &2
 \end{array}
\]
so $g_0=1$ and every later gap is even, the primes after $2$ being odd.
A second normalisation with denominator $2^{\,i}$ also occurs,
and at every finite horizon it is exactly twice this one,
\[
 \sum_{i=0}^{n-1}\frac{p_i}{2^{\,i}}
 =2\sum_{i=0}^{n-1}\frac{p_i}{2^{\,i+1}} ,
\]
the
\href{https://github.com/wcook04/plectis-erdos/blob/3d6d938d696fed0fb71dd55115a18a73738ff223/lean/ErdosProblems/Erdos251/PrimeGapDyadicTail.lean\#L111}{factor-of-two
normalisation}.  A factor of two does not change rationality, so no result
below depends on the choice; the identity is stated so that the indexing cannot
drift silently.

\paragraph{Notation.}
$\N=\{0,1,2,\ldots\}$ and $\Npos=\{1,2,\ldots\}$.  We write $\ph$ for Euler's
totient function, $\operatorname{den}q$ for the denominator of a rational
number $q$ written in lowest terms, $\operatorname{dist}(x,\Z)$ for the
distance from a real number $x$ to the nearest integer, and
$\operatorname{Irr}(x)$ for the assertion that $x$ is irrational.  A rational
number is called \emph{integral} when it is the image of an integer.  A
sequence $(a_n)$ is \emph{eventually periodic with period $h\ge1$} when
$a_{n+h}=a_n$ for every sufficiently large $n$.  A set $A\subseteq\N$ has
\emph{density zero} when $|A\cap[1,x]|=o(x)$, and a property holds for
\emph{almost all} indices when its exceptional set has density zero.  A sum
over an empty range of indices is zero.

\section{Sparse congruence-preserving perturbations}
\label{long251:sec:paired-corrections}

The sparse construction preserves cumulative residues as well as coefficient
residues.  Two adjacent corrections provide a free weighted digit while
fixing their unweighted total.  A third coordinate is needed to repair the
cumulative residue before that pair.  We account separately for its location,
size and contribution to the series.

\paragraph{Density and block conventions.}
Here $\N=\{0,1,\ldots\}$, all intervals of indices contain integers, and
$\log$ is natural unless a base is displayed.  Upper Banach density zero means
\[
 \lim_{H\to\infty}\sup_{u\in\N}\frac{|S\cap[u,u+H)|}{H}=0.
\]
For integers $X\ge1$ and $m\ge1$, let $\mu_{a,X,m}$ be the empirical
probability measure giving mass $1/X$ to each literal block
$(a_n,\ldots,a_{n+m-1})$, $X\le n<2X$, with multiplicities.
We use $d_{\rm TV}(\mu,\nu)=\sup_B|\mu(B)-\nu(B)|$; consequently tests
bounded in absolute value by $B_0$ differ in mean by at most
$2B_0d_{\rm TV}(\mu,\nu)$.  No rescaling of the coefficients is implicit.

\begin{theorem}[sparse rationalisation and congruences]
\label{long251:res:sparse-rationalisation}
Let $a_n\in\N$ and $A=\sum_{n\ge0}a_n2^{-n-1}<\infty$.  For every cutoff
$K$ and every $f:\N\to\R$ tending to infinity, there are a set
$S\subseteq[K,\infty)$ of upper Banach density zero and a nondegenerate
interval $I\subset(A,\infty)$ with the following property.  Every $r\in I$
has the form
\[
 r=\sum_{n\ge0}(a_n+e_n)2^{-n-1},
\]
where $e_n\in\N$ is supported on $S$, $e_n\le f(n)$ eventually, and, for
every integer $q\ge1$, there is an $N_q$, chosen independently of $r$, such that
\[
 q\mid e_n\quad\hbox{and}\quad q\mid\sum_{i<n}e_i
 \qquad(n\ge N_q).
\]
For every $\varepsilon>0$, the construction can instead be chosen with
$e_n\le(\log(n+3))^\varepsilon$ eventually and
$|S\cap[X,2X)|=O_\varepsilon(X/\log\log X)$.  In that case the empirical
distributions of original and corrected unnormalised blocks of length
$m(X)=o(\log\log X)$, sampled at the same integer starts in $[X,2X)$,
have total variation distance tending to zero, uniformly over $r$.
Precisely the coupling error is at most $m|S\cap[X,2X+m)|/X$; with
$\|\Phi\|_\infty\le B_0$ the bounded-test error is at most twice this
quantity times $B_0$, also for tests depending on the starting index.
\end{theorem}

\noindent The supplied index marks the general and polylogarithmic
existential endpoints as \texttt{ci\_checked}; precise declarations follow
the proof.  The printed buffered-triple witness and its schedule are an
ordinary reconstruction, not the witness used by that formal development.
The statement correspondence includes the target-independent support and
the growing-block quantifiers; it does not identify the constructions.

\begin{proof}
Choose increasing centres $n_j$, indexed by $j\ge0$, together with an initial
$n_{-1}\ge K$.  Their spacings $s_j=n_j-n_{j-1}$ will be at least $4$.
Let $M_j$ be positive integer moduli with $M_j\mid M_{j+1}$, and set
$D_j=2^{s_j}-1$.  These parameters will be specified below independently of
all free digits.

\emph{Place the buffer before the pair.}
Set $C_0=0$ and recursively define the deterministic integers
\[
 c_j=(-C_j)\bmod M_j,\qquad C_{j+1}=C_j+c_j+M_jD_j.
\]
For any choice $d_j\in\{0,\ldots,D_j\}$, put
\begin{equation}\label{long251:eq:sparse-triple}
 e_{n_j-1}=c_j,\qquad e_{n_j}=M_jd_j,\qquad
 e_{n_j+1}=M_j(D_j-d_j),
\end{equation}
and put $e_n=0$ off these triples.  The spacing makes the triples disjoint.
Their total is $c_j+M_jD_j$, independent of $d_j$, so $C_j$ really is the
cumulative correction $C_j=\sum_{i<n_j-1}e_i$ before its buffer, for every choice of free digits.
The pair alone could not repair that cumulative residue: its total is
already a multiple of $M_j$.  The buffer does, since $M_j\mid C_j+c_j$.

For example, $M=2,D=3$ permits the pairs $(0,6),(2,4),(4,2),(6,0)$.
Each has total $6$, but their weighted contributions have numerators
$6,8,10,12$ over $2^{n+2}$.  A preceding cumulative total $C=1$ needs the
additional buffer $c=1$ at $n-1$; its weighted contribution is $2^{-n}$.
This fixed contribution changes the location of the attainable interval,
not its free-digit capacity.

\emph{Check every late coordinate.}
Assume that each fixed $q$ divides $M_j$ eventually, and choose $J$ with
$q\mid M_J$.  At index $n_J$, just after its buffer, the cumulative sum
is divisible by $q$.  Both entries of the pair are also divisible by $q$,
so every later cumulative sum through that stage remains divisible by $q$.
Inductively $q\mid C_j$ and $q\mid M_j$ imply
$q\mid c_j$, since $M_j\mid C_j+c_j$.  Thus every subsequent buffer,
pair entry and zero entry is divisible by $q$, as is every intermediate
cumulative sum.  We may take $N_q=n_J$.  The first buffer need not itself
be divisible by $q$, and its coordinate $n_J-1$ is deliberately outside
this eventual range.

\emph{Choose a schedule for arbitrary $f$.}
Define
\[
 h(n)=\inf_{m\ge n}\min(f(m),m).
\]
This finite real number is nondecreasing in $n$, satisfies $h(n)\le n$,
and tends to infinity.  Choose $n_{-1}\ge K$ with $h(n_{-1})\ge32$ and
recursively set
\begin{equation}\label{long251:eq:sparse-general-schedule}
 \begin{split}
 k_j&=\max\{k\in\N:k\ge2,\ k!2^{k+2}\le h(n_{j-1})\},\\
 M_j&=k_j!,\qquad s_j=k_j+2,\qquad n_j=n_{j-1}+s_j.
 \end{split}
\end{equation}
The maximum exists: $k=2$ is eligible and $k!2^{k+2}$ tends to infinity.
Since $h$ is nondecreasing, the integers $k_j$ are nondecreasing; since
$n_j\to\infty$, they tend to infinity.  The factorials form a divisibility
chain and eventually contain every fixed divisor, as do the distinguished
terms in~\cite[Proposition~8]{vandoornkovac2025}.  At each of the three
coordinates $n$ in stage $j$,
\[
 0\le e_n\le M_j2^{s_j}\le h(n_{j-1})\le f(n),
 \qquad e_n\le n_{j-1}<n.
\]
Here the buffer is smaller than $M_j$ and each pair entry is at most
$M_jD_j$.  Thus every correction series converges by comparison with
$\sum_n n2^{-n-1}$.

Let $S=\bigcup_j\{n_j-1,n_j,n_j+1\}$.  It is independent of the digits.
To see upper Banach density zero, fix $H>0$.  Outside a finite initial
segment, consecutive centres are at least $H$ apart, since $s_j\to\infty$.
An interval of length $L$ therefore meets at most
$3(L/H+2)$ late support coordinates, plus the fixed finite number of early
coordinates.  Divide by $L$, take the supremum over interval positions,
then let $L\to\infty$ and $H\to\infty$.

\emph{Retain an interval of values.}
Put $w_j=M_j2^{-n_j-2}$ and
\[
 F_j=\sum_{i\ge j}D_iw_i,\qquad
 \beta=\sum_{j\ge0}\bigl(c_j2^{-n_j}+D_jw_j\bigr).
\]
The size estimates just proved give convergence of both series, $F_j\to0$,
and $F_0>0$.  The baseline $\beta$ is positive and independent of the digits.
The complete weighted correction, including every buffer, is exactly
\begin{equation}\label{long251:eq:sparse-baseline}
 \sum_n e_n2^{-n-1}=\beta+\sum_jd_jw_j.
\end{equation}
For $i>j$, we have $M_i\ge M_j$ and
\[
 D_iw_i=M_i\bigl(2^{-n_{i-1}-2}-2^{-n_i-2}\bigr)
 \ge M_j\bigl(2^{-n_{i-1}-2}-2^{-n_i-2}\bigr).
\]
Telescoping over $i=j+1,\ldots,J$ and letting $J\to\infty$ yields
$F_{j+1}\ge w_j$.  This is the needed overlap condition.  The intervals
$[dw_j,dw_j+F_{j+1}]$, for $0\le d\le D_j$, cover $[0,F_j]$.
For any $y\in[0,F_0]$, choose successive digits leaving each remainder in
$[0,F_{j+1}]$.  Since these capacities tend to zero, the resulting series
has sum $y$.  In the notation of Crmari\'{c}--Kova\v{c}
\cite[Lemma~4]{crmarickovac2025}, the finite menu
$X_j=\{0,w_j,\ldots,D_jw_j\}$ has largest successive gap $w_j$ and
diameter $D_jw_j$, so its covering condition is precisely
$w_j\le F_{j+1}$.  Fridy's generalised-base lemma
\cite[Lemma, p.~194]{fridy1966} and the reciprocal-menu argument in
Kova\v{c}--Tao \cite[Lemma~5.1]{kovactao2024} are antecedents.
No monotonicity of the weights is needed for the finite-menu proof.
Equation~\eqref{long251:eq:sparse-baseline} therefore realises every $r$ in
the fixed interval $I=(A+\beta,A+\beta+F_0)\subset(A,\infty)$.

\emph{Choose a polylogarithmic schedule.}
Fix $0<\varepsilon<1$.  Replace~\eqref{long251:eq:sparse-general-schedule} by
\begin{equation}\label{long251:eq:sparse-polylog-schedule}
 \begin{split}
 s_j&=\left\lfloor\frac{\varepsilon}{2}
          \log_2\log(n_{j-1}+3)\right\rfloor,\qquad n_j=n_{j-1}+s_j,\\
 k_j&=\max\{k\in\N:k\ge2,\ k!\le(\log(n_{j-1}+3))^{\varepsilon/4}\},
 \qquad M_j=k_j!.
 \end{split}
\end{equation}
Take the initial cutoff so large that $s_j\ge4$ and the maximum is nonempty
from the first step.  Again $k_j$ is nondecreasing and tends to infinity.
Keep $D_j=2^{s_j}-1$ and the same buffer recursion.  All triple entries obey
\[
 e_n\le M_j2^{s_j}\le(\log(n_{j-1}+3))^{3\varepsilon/4}
 \le(\log(n+3))^\varepsilon.
\]
These polylogarithmic bounds still give convergence of $F_j$ and $\beta$.
Every argument above concerning congruences and interval filling remains
valid.  Also $s_j$ is comparable, with constants depending on $\varepsilon$,
to $\log\log(n_{j-1}+3)$, and $s_j=o(n_{j-1})$.  Hence consecutive centres
near $X$ are separated by at least a constant multiple of
$\log\log X$, giving $|S\cap[X,2X)|=O_\varepsilon(X/\log\log X)$.
The same estimate holds on $[X,3X)$.

Finally sample both length-$m$ blocks at the same uniform integer start in
$[X,2X)$.  They can differ only if that block meets $S$.  Each changed
coordinate belongs to at most $m$ such blocks, so the total variation
distance is at most
\[
 \frac{m\,|S\cap[X,2X+m)|}{X}
 =O_\varepsilon\!\left(\frac{m}{\log\log X}\right).
\]
For integer $m=m(X)=o(\log\log X)$ this tends to zero.  This comparison
uses the literal, unnormalised blocks and requires no limiting distribution
for the original sequence.
\end{proof}

The pair identities, modular repair and conditional interval filling have
formal counterparts in \path{SparseRationalisationCore.lean}.  The supplied
index marks \href{https://github.com/wcook04/plectis-erdos/blob/3d6d938d696fed0fb71dd55115a18a73738ff223/lean/ErdosProblems/Erdos251/SparseAmbientR9.lean\#L217}{the arbitrary-word interval theorem}, \href{https://github.com/wcook04/plectis-erdos/blob/3d6d938d696fed0fb71dd55115a18a73738ff223/lean/ErdosProblems/Erdos251/SparsePaperR11.lean\#L53}{its rational-target consequence}, \href{https://github.com/wcook04/plectis-erdos/blob/3d6d938d696fed0fb71dd55115a18a73738ff223/lean/ErdosProblems/Erdos251/SparsePaperR11.lean\#L78}{the polylogarithmic interval theorem}, \href{https://github.com/wcook04/plectis-erdos/blob/3d6d938d696fed0fb71dd55115a18a73738ff223/lean/ErdosProblems/Erdos251/GrowingBlocksR11.lean\#L139}{the block total-variation bound} and \href{https://github.com/wcook04/plectis-erdos/blob/3d6d938d696fed0fb71dd55115a18a73738ff223/lean/ErdosProblems/Erdos251/GrowingBlocksR11.lean\#L152}{the uniform bounded-test theorem} as \texttt{ci\_checked}.
The full existential conclusions are therefore covered by the supplied
verification record.  The formal construction is a one-site
residue-feedback construction, not the buffered-triple schedule printed
above.  Agreement of theorem statements must not be confused with identity
of these witnesses.  There is no Comparator entry for this theorem.
Applied to
prime gaps, the construction makes no assertion that the corrected
cumulative positions remain prime.

\begin{corollary}[arbitrarily local target intervals]\label{long251:res:local-targets}
For the arbitrary-word sparse theorem, the target interval can additionally
be required to lie in $(A,A+\eta)$ for any prescribed $\eta>0$.
\end{corollary}
\begin{proof}
In the general schedule, every supported entry satisfies $e_n\le n$.
Start beyond an integer $K'$ with
$\sum_{n\ge K'}n2^{-(n+1)}=(K'+1)2^{-K'}<\eta$.
The complete interval of corrections has positive lower endpoint and
upper endpoint at most this sum.  This strengthens the location of the
interval, not a claim that every number above $A$ is attainable under
one fixed envelope.
\end{proof}

\section{Context and mathematical dependencies}\label{long251:sec:context}

\paragraph{Relation to prior work.}
The passage from the primes to their gaps is already public.  Tao posted on
the problem's forum thread on 7 October 2025 that summation by parts makes the
question equivalent to irrationality of $\sum_n(p_{n+1}-p_n)2^{-n}$, and named
the shape of the missing input as a sufficiently quantitative and uniform
prime-tuples hypothesis giving statistical control of the binary expansion of
about $\log\log n$ consecutive
gaps~\cite[comment of 7 October 2025]{erdosproblems251thread}.
Theorem~\ref{long251:res:infinite} below is that reduction in exact form, with the
endpoint retained, with convergence supplied by an elementary bound, and with
the statement checked by the Lean kernel.  The elementary bound replaces the
prime number theorem in this reduction.

A sufficiently uniform Hardy--Littlewood hypothesis gives conditional
results of a different kind.  Land's draft states conditional
irrationality \cite[Conjecture~1, Theorem~2]{land2026}; Ringer's draft states
conditional normality in each integer base \cite{ringer2026}.  The authors
supply formalisation material, which has not been independently rebuilt
for this revision.  The uniform hypothesis in
\cite[Conjecture~1.3]{kuperberg2023} controls growing tuples, with constants
chosen before the tuple and basepoint; qualitative fixed-tuple
Hardy--Littlewood does not provide it.

The version of Ringer's source consulted on 16 September 2026 separates a
stopped first-point comparison from this stronger conjecture.  Its averaged
one-sided tuple error is sufficient, and a bounded-domination condition on
first-point masses is weaker still; these are alternative sufficient
inputs, not conclusions of presently invoked unconditional prime-gap
results.  The signs matter: both an undercount at one parity of the
truncation and an overcount at the other can be adverse.  An upper sieve
alone therefore does not supply that averaged input.  None of these
conditional inputs is assumed below.

Erd\H{o}s stated on p.~93 of the 1958 article that $\sum_n p_n^{k}/n!$ is
irrational for every $k\ge1$, and wrote there that the proof for $k>1$ is
complicated enough that only the case $k=1$ is
printed~\cite[pp.~93--95]{erdos1958}.  A proof of the full family appears in
Schlage-Puchta's Theorem~3, which gives the stronger statement that
$1,S_0,S_1,S_2,\ldots$ are linearly independent over $\Q$, where
$S_k=\sum_{n\ge1}p_n^{k}/n!$~\cite[Theorem~3]{schlagepuchta2011}.

On page~103 of his 1988 problem paper Erd\H{o}s separately stated the
fixed-denominator problem: he could not prove that $\sum_n p_n^{k}/2^{n}$ is
irrational for every $k\ge1$, and wrote that the case $k=1$ was probably
already very difficult.  He also stated the variable-denominator expectation
that $\sum_{n\ge1}p_n/(g_1\cdots g_n)$ is irrational whenever $g_n\ge2$ and
$g_n=o(p_n)$~\cite[p.~103]{erdos1988}.  That expectation is false: on
15 April 2026 Kova\v{c} posted a note, with the printed attribution
ChatGPT~5.4 Pro (orchestrated by Vjeko Kova\v{c}), constructing such a
sequence $(g_n)$ for which the sum is exactly
$1$~\cite[Theorem~1 and proof, pp.~1--2]{kovac2026}.

Section~3 of the 1958 article concerns variable product denominators,
not the dyadic denominator sequence \cite[pp.~96--97]{erdos1958}.
Its printed growth condition (5) uses nonstandard asymptotic typography;
we do not use that condition as an operative hypothesis or silently
translate it into a modern quantified statement.  The simple endpoint
example is unambiguous: if $G_n=\prod_{j=1}^n(p_j+1)$ and $G_0=1$, then
\[
 \frac{p_n}{G_n}=\frac1{G_{n-1}}-\frac1{G_n},\qquad
 \sum_{n\ge1}\frac{p_n}{G_n}=1.
\]
Neither this variable-denominator example nor the counterexample in
\cite{kovac2026} addresses fixed dyadic denominators.

There is also a genuinely adjacent proved dyadic theorem.  If $P(m)$ denotes
the largest prime factor of $m$, Erd\H{o}s and Pomerance proved that
\[
 \sum_{m\ge2}\frac{\mathbf 1_{\{P(m)>P(m+1)\}}}{2^m}
\]
is irrational~\cite[\S7, p.~320]{erdospomerance1978}.  Erd\H{o}s and Graham
record the complementary indicator on p.~62: equality of the two largest prime
factors is impossible for consecutive integers, so its series is $1/2$ minus
the displayed one and is irrational as well.  That is a theorem about a bounded
prime-factor comparison digit sequence, and the numerators in $\Pi$ are
unbounded.

Two further results of Schlage-Puchta enter below.  His Theorem~2 classifies
rationality for numbers whose base-$b$ digit string concatenates the
representations of a slowly growing integer sequence, and is a different
object from the tail criterion here; it was already pointed at this problem in
the catalogue thread on 15 April 2026.  His Lemma~4 is the input this note
actually uses.  It states that for every polynomial $F\in\Z[x_0,\ldots,x_k]$
which does not vanish identically, $F(g_n,\ldots,g_{n+k})\ne0$ for almost all
$n$, where $g_n$ are the actual consecutive prime
gaps~\cite[Lemma~4, pp.~5--6]{schlagepuchta2011}.  Its proof runs through
Selberg's sieve.  Sections~\ref{long251:sec:obstructions} draws two consequences from
it.

Theorems~\ref{long251:res:boundedperturbation}
and~\ref{long251:res:sparse-rationalisation} fill intervals of attainable
values.  The first is a binary expansion.  For the second, the direct
reference is Crmari\'{c}--Kova\v{c}'s finite-choice covering lemma
\cite[Lemma~4]{crmarickovac2025}: a menu's largest successive gap must not
exceed the total remaining diameter.  They apply it to product-denominator
series.  Our menus are $X_j=\{0,w_j,\ldots,D_jw_j\}$, so the hypothesis is
$w_j\le\sum_{i>j}D_iw_i$.  Fridy's generalised-base lemma
\cite[Lemma, p.~194]{fridy1966} is an antecedent with nonincreasing weights;
the finite-menu form avoids that additional assumption.  Kova\v{c}--Tao use
reciprocal menus \cite[Lemma~5.1, Theorem~2.5]{kovactao2024}, while van
Doorn--Kova\v{c} combine filling with a divisibility chain absorbing every
integer \cite[Lemma~7, Proposition~8]{vandoornkovac2025}.  Our factorial
moduli use the latter absorption property to preserve eventual congruences.

A further comparison is van Doorn's exchange
$\{21d,28d\}\leftrightarrow\{20d,30d\}$: both reciprocal sums are $1/(12d)$,
while the ordinary sums are $49d$ and $50d$
\cite[proof of Theorem~3]{vandoorn2025}.  Our paired corrections preserve
the ordinary sum and vary the dyadic contribution instead.  This is an
analogy between conserved quantities, not an application of his partition
theorem.  We claim no new interval-covering principle.

Binary achievement sets require care beyond these covering inequalities.
Bartoszewicz--Filipczak--Szymonik use a central run of attainable integer
block sums to obtain interior in a multigeometric family
\cite[Theorem~2]{bartoszewicz2014}.  Prus-Wi\'sniowski--Ptak construct
Cantorvals for which the overlap indices
$\{n:a_n\le\sum_{i>n}a_i\}$ have density zero \cite[Theorem]{prusptak2024}.
Thus neither occasional separated cylinders nor failure of an eventual
term-versus-tail alternative alone decides interior.  The survey
\cite[\S1]{glabprus2025} records the surrounding classification.  Our
variable-digit proof verifies covering at every stage and does not depend
on a classification theorem for binary subsums.  Nitecki's preprint
\cite[Theorem~14]{nitecki2015} gives an exposition explicitly crediting the
Guthrie--Nymann classification; its title and numbering differ from the
published Monthly article.

The inequality pattern alone is insufficient even when its index set is
specified exactly: Marchwicki--Miska's construction
\cite[Theorem~2.1]{marchwickimiska2021}, with the repaired proof in
\cite{miskaprusptak2023}, can realise any infinite strict term-dominating
index set with a Cantor achievement set.  The repair preserves the original
uniqueness conclusion; the latter paper also gives a simpler proof of a
weaker statement without uniqueness.  Nowakowski's Star Procedure
\cite[Definition~2, Theorem~3.1]{nowakowski2025} is a different sufficient
route to a Cantorval, requiring every stage of a positivity recursion.
None of these results is asserted to settle the prescribed factorial
binary family in the separate working note.

\paragraph{Automatic sequences and the limit of the analogy.}
Adamczewski--Drmota--M\"ullner prove computable logarithmic densities for
fixed automatic sequences along primes, and give criteria for natural
densities and their rationality \cite[Theorems~1.2, 1.4]{adm2022}.
The general logarithmic densities are not all rational.  The fixed
finite-alphabet automaton is essential: the result is not a theorem about
unbounded consecutive prime gaps, their complete dyadic tails, or a growing
family of residue or carry automata.  Any such application would have to
specify the automaton, prove that it computes the required observable, and
supply uniform errors as its state space grows.  We do not infer any of
those steps from the density theorem.

\paragraph{The uniformity required from prime statistics.}
Kuperberg's large-set singular-series estimates allow a specified growing
cardinality, not arbitrary dimension \cite[Theorem~1.1]{kuperberg2023}.
Her arithmetic-progression and smooth-weight estimates are a useful
fixed-parameter comparison \cite{kuperbergweighted2025}; they do not by
themselves give a law for ordered consecutive-gap blocks.  Jha's revised
Poisson-tail preprint \cite{jha2026} concerns growing short-interval prime
counts under a strong Hardy--Littlewood hypothesis.  Counting primes in an
interval and controlling the joint position-weighted gap-tail observable
are different tasks.  No growing-window estimate needed below is claimed
to follow merely from these related titles or from a fixed-dimensional
limit theorem.

Problem~\#251 also appears as the unproved declaration \texttt{erdos\_251} in
the \emph{Formal Conjectures} repository~\cite{formalconjectures251}.  Its
zero-based Lean sum starts with the zeroth prime over $2^0$, so it is twice the
displayed normalisation and has equivalent irrationality status; its proof is
\texttt{sorry}.

\paragraph{The strategy.}
The digits $p_n$ grow, so the series is not a digit expansion in any bounded
alphabet, and the standard rationality criteria for such expansions do not
apply directly.  The classical elementary criteria for series of this kind
instead control irrationality through the growth of the denominators.
Erd\H{o}s and Straus named a sum $\sum_k 1/a_k$ over a strictly increasing
sequence of positive integers an \emph{Ahmes series}~\cite{erdosstraus1963};
for such a series the condition $a_k^{1/2^k}\to\infty$ is sufficient for
irrationality, and it is sharp, since shifted Sylvester sequences grow like
$C^{2^{k}}$ for arbitrarily large $C$ and have rational reciprocal sum.  Both
statements and their attribution are recorded in the introduction of Kova\v{c}
and Tao~\cite[\S1]{kovactao2024}.  Splitting each term $p_i/2^{\,i+1}$ into
$p_i$ copies of $2^{-(i+1)}$ writes $\Pi$ as a sum of unit fractions with
repetitions, and the denominators occurring in it are exactly the powers of
two.  Even ignoring the repetitions the growth hypothesis fails by every
available margin, since $(2^{\,n})^{1/2^{\,n}}\to1$, and every arithmetic
constraint has to come from the numerators instead.

What replaces growth control is denominator control on the sequence of
rescaled tails.  Rationality of the sum is equivalent to an eventual
integrality condition on differences of those tails, and that condition uses
nothing about the numerators beyond the fact that they are integers.  The
condition does not by itself force the numerators to repeat:
Proposition~\ref{long251:res:telescope}, applied to $K_n=n$, produces the integer
sequence $\kappa_n=n-1$, which is unbounded and hence not eventually periodic,
and whose dyadic sum is zero.

\paragraph{Dependencies and reading routes.}
The sparse theorem is independent of primes: modular repair and finite-menu
covering give the interval, and counting changed starts gives its statistical
consequences.  Schlage-Puchta's lemma and the prime number theorem enter
only the prime-anchored joint corollary.
The recurrence classification needs integer digits and Euler's congruence,
not any prime-tuple conjecture.  The one-tail certificate adds an explicit
majorant and exact finite prime data.  Irrationality would require the
additional cofinal prime-gap input in Section~\ref{long251:sec:open}.
The appendices give the finite certificates, complete moved statements and
source concordance.  The accompanying statement-destination table preserves
the original short-note labels at these locations.

\section{Summation by parts, with the endpoint retained}\label{long251:sec:parts}

Summation by parts trades a sequence for its consecutive differences.  The form
recorded here is exact: it carries no error term, it assumes nothing about the
sequence, and it retains the endpoint term.  No endpoint is absorbed into an
estimate.  For a sequence $P$ of rational numbers and $n\ge0$ put
\[
 D(P,n)=\sum_{i=0}^{n-1}\frac{P(i)}{2^{\,i+1}},
 \qquad
 \Delta(P,n)=\sum_{i=0}^{n-1}\frac{P(i+1)-P(i)}{2^{\,i+1}} ,
\]
both empty, hence zero, at $n=0$.  We call these the
\href{https://github.com/wcook04/plectis-erdos/blob/3d6d938d696fed0fb71dd55115a18a73738ff223/lean/ErdosProblems/Erdos251/PrimeGapDyadicTail.lean\#L101}{dyadic partial
sum} and the
\href{https://github.com/wcook04/plectis-erdos/blob/3d6d938d696fed0fb71dd55115a18a73738ff223/lean/ErdosProblems/Erdos251/PrimeGapDyadicTail.lean\#L121}{dyadic
difference sum} of $P$.

\begin{proposition}[finite summation by parts]\label{long251:res:abel}\compatlabel{res:abel}
For every rational sequence $P$ and every $n\ge0$,
\[
 D(P,n+1)=P(0)+\Delta(P,n)-\frac{P(n)}{2^{\,n+1}} .
\]
\end{proposition}

\begin{proof}
A routine induction on $n$.  At $n=0$ both sides equal $P(0)/2$.  For the step,
adding $P(n+1)/2^{\,n+2}$ to the left and $\bigl(P(n+1)-P(n)\bigr)/2^{\,n+1}$
to the difference sum changes the endpoint term from $P(n)/2^{\,n+1}$ to
$P(n+1)/2^{\,n+2}$, and the two adjustments agree.
\end{proof}

\noindent Formalised as the
\href{https://github.com/wcook04/plectis-erdos/blob/3d6d938d696fed0fb71dd55115a18a73738ff223/lean/ErdosProblems/Erdos251/PrimeGapDyadicTail.lean\#L138}{summation-by-parts
identity}.  Nothing is assumed about $P$: no positivity, no monotonicity, and
no convergence.

Specialising to $P(i)=p_i$, whose first value is $p_0=2$, and writing
$g_i=p_{i+1}-p_i$ for the zero-based gaps, gives the reformulation.

\begin{theorem}[prime-gap reformulation]\label{long251:res:parts}\compatlabel{res:parts}
Let $p_0=2,p_1=3,\ldots$ be the primes in increasing order and
$g_i=p_{i+1}-p_i$.  For every $n\ge0$,
\[
 \sum_{i=0}^{n}\frac{p_i}{2^{\,i+1}}
 =2+\sum_{i=0}^{n-1}\frac{g_i}{2^{\,i+1}}-\frac{p_n}{2^{\,n+1}} .
\]
\end{theorem}

\noindent Formalised as the
\href{https://github.com/wcook04/plectis-erdos/blob/3d6d938d696fed0fb71dd55115a18a73738ff223/lean/ErdosProblems/Erdos251/PrimeGapDyadicTail.lean\#L172}{prime-gap
summation by parts}, using the
\href{https://github.com/wcook04/plectis-erdos/blob/3d6d938d696fed0fb71dd55115a18a73738ff223/lean/ErdosProblems/Erdos251/PrimeGapDyadicTail.lean\#L161}{gap partial
sum} and the
\href{https://github.com/wcook04/plectis-erdos/blob/3d6d938d696fed0fb71dd55115a18a73738ff223/lean/ErdosProblems/Erdos251/PrimeGapDyadicTail.lean\#L156}{gap cast
identity}; the latter records that the natural-number difference
$p_{n+1}-p_n$ agrees with the difference taken in $\Q$, which needs
$p_n\le p_{n+1}$, the
\href{https://github.com/wcook04/plectis-erdos/blob/3d6d938d696fed0fb71dd55115a18a73738ff223/lean/ErdosProblems/Erdos251/PrimeGapDyadicTail.lean\#L152}{monotonicity
step}.  The leading $2$ is the first prime, not a normalising constant.
At $n=2$, for instance, the left side is $2/2+3/4+5/8=19/8$ and the right side
is $2+(1/2+2/4)-5/8=19/8$.

\subsection{The infinite identity and the irrationality equivalence}\label{long251:sec:infinite}\compatlabel{sec:infinite}

Write
\[
 u_n=\frac{p_n}{2^{\,n+1}},\qquad
 v_n=\frac{g_n}{2^{\,n+1}}
\]
for the terms of the prime series and of the gap series, so that
$\sum_{n\ge0}u_n=\Pi$.  The termwise identity $v_n=2u_{n+1}-u_n$ is the
\href{https://github.com/wcook04/plectis-erdos/blob/3d6d938d696fed0fb71dd55115a18a73738ff223/lean/ErdosProblems/Erdos251/PrimeGapDyadicTail.lean\#L202}{dyadic
discrete derivative}.  It expresses each gap term as an integer combination of
two consecutive prime terms, so once $(u_n)$ is summable the gap series can be
summed by rearranging two copies of the prime series.

\begin{theorem}[infinite prime-gap identity]\label{long251:res:infinite}
Both series converge and
\[
 \sum_{n\ge0}\frac{p_n}{2^{\,n+1}}
 \;=\;2+\sum_{n\ge0}\frac{g_n}{2^{\,n+1}} .
\]
\end{theorem}

\begin{proof}
The bound $p_n\le1250(n+1)^4$ of Appendix~\ref{long251:app:prime-bound} makes $(u_n)$
summable, and $0<g_n\le p_{n+1}$ then makes $(v_n)$ summable.  The shifted
sequence $(u_{n+1})$ is summable, so summing $v_n=2u_{n+1}-u_n$ and using
$u_0=1$ gives
\[
 \sum_{n\ge0}v_n
 =2\left(\sum_{n\ge0}u_n-u_0\right)-\sum_{n\ge0}u_n
 =\sum_{n\ge0}u_n-2 . \qedhere
\]
\end{proof}

\noindent The polynomial bound is
\href{https://github.com/wcook04/plectis-erdos/blob/3d6d938d696fed0fb71dd55115a18a73738ff223/lean/ErdosProblems/Erdos251/PrimeGapDyadicTail.lean\#L360}{checked
here}, summability of the prime series
\href{https://github.com/wcook04/plectis-erdos/blob/3d6d938d696fed0fb71dd55115a18a73738ff223/lean/ErdosProblems/Erdos251/PrimeGapDyadicTail.lean\#L379}{here},
the summability transfer is the
\href{https://github.com/wcook04/plectis-erdos/blob/3d6d938d696fed0fb71dd55115a18a73738ff223/lean/ErdosProblems/Erdos251/PrimeGapDyadicTail.lean\#L385}{gap-series
summability theorem}, and the displayed identity is the
\href{https://github.com/wcook04/plectis-erdos/blob/3d6d938d696fed0fb71dd55115a18a73738ff223/lean/ErdosProblems/Erdos251/PrimeGapDyadicTail.lean\#L427}{unconditional
prime-gap identity}.  The prime number theorem is needed neither for this
convergence nor for the identity; it enters later through $p_n\sim n\log n$
in Corollary~\ref{long251:res:nonconc-primes} and in the state-compression
argument~\cite[Chapter~6]{mv2007}.

\begin{corollary}[exact irrationality reformulation]\label{long251:res:irr-equivalence}\label{res:irr-equivalence}
$\Pi$ is irrational if and only if $S=\sum_{n\ge0}g_n2^{-(n+1)}$ is irrational.
The corresponding zero-based series with denominator $2^n$ equals $4+2S$ and
has the same irrationality status.
\end{corollary}

\noindent These are the
\href{https://github.com/wcook04/plectis-erdos/blob/3d6d938d696fed0fb71dd55115a18a73738ff223/lean/ErdosProblems/Erdos251/PrimeGapDyadicTail.lean\#L435}{normalised
irrationality equivalence}, the
\href{https://github.com/wcook04/plectis-erdos/blob/3d6d938d696fed0fb71dd55115a18a73738ff223/lean/ErdosProblems/Erdos251/PrimeGapDyadicTail.lean\#L444}{displayed-series
identity}, and the
\href{https://github.com/wcook04/plectis-erdos/blob/3d6d938d696fed0fb71dd55115a18a73738ff223/lean/ErdosProblems/Erdos251/PrimeGapDyadicTail.lean\#L459}{displayed-series
irrationality equivalence}.  Concretely, $\Pi=3.674643966\ldots$ is irrational
if and only if $S=1.674643966\ldots$ is, and neither is known.

\section{The tail recurrence and the exact criteria}\label{long251:sec:tail}

The series is not attacked directly.  Suppose $\sum_{i\ge0}a_i2^{-(i+1)}$
converges with every $a_i$ an integer, and rescale its tails by putting
\[
 T_N=2^{\,N+1}\sum_{i>N}\frac{a_i}{2^{\,i+1}}
    =\sum_{j\ge1}\frac{a_{N+j}}{2^{\,j}} .
\]
Two facts follow immediately.  First $T_{N+1}=2T_N-a_{N+1}$, so moving one
level along doubles the rescaled tail and subtracts a single coefficient.
Second $T_0=2\sum_{i\ge0}a_i2^{-(i+1)}-a_0$, so the sum is rational exactly
when $T_0$ is.  This section studies that recurrence, assuming nothing about
the coefficients beyond the fact that they are integers, and resumes the
prime-gap instance at the end.

\begin{definition}\label{long251:def:rec}\compatlabel{def:rec}
Let $g:\N\to\Z$ and $T:\N\to\Q$ or $T:\N\to\R$.  Say $T$ satisfies the
\emph{dyadic tail recurrence} with \emph{digits} $g$ when
$T_{N+1}=2T_N-g_{N+1}$ for every $N$.  Write $\sigma_h(N)=T_{N+h}-T_N$ for the
\emph{shift} of length $h$ at $N$, and call a rational number \emph{integral}
when it is the image of an integer.
\end{definition}

\noindent These are the
\href{https://github.com/wcook04/plectis-erdos/blob/3d6d938d696fed0fb71dd55115a18a73738ff223/lean/ErdosProblems/Erdos251/PrimeGapDyadicTail.lean\#L482}{tail
recurrence}, the
\href{https://github.com/wcook04/plectis-erdos/blob/3d6d938d696fed0fb71dd55115a18a73738ff223/lean/ErdosProblems/Erdos251/PrimeGapDyadicTail.lean\#L544}{shift}, and
\href{https://github.com/wcook04/plectis-erdos/blob/3d6d938d696fed0fb71dd55115a18a73738ff223/lean/ErdosProblems/Erdos251/PrimeGapDyadicTail.lean\#L575}{integrality}.
One small orbit is worth having in view.  Take every digit $g_N=0$ and
$T_0=1/12$.  Then the orbit is
$\tfrac1{12},\tfrac16,\tfrac13,\tfrac23,\tfrac43,\ldots$ and
$\sigma_h(N)=(2^{h}-1)2^{\,N}/12$.  The shift of length $2$ is not integral at
$N=0$ and is integral at $N=2$; the shift of length $1$ is never integral.
Nonintegrality at one prescribed shift length therefore does not imply
irrationality.

Iterating the recurrence $h$ times multiplies $T_N$ by $2^{h}$ and accumulates
an explicit integer.  Define $B_{0,N}=0$ and $B_{h+1,N}=2B_{h,N}+g_{N+h+1}$,
so that $B_{h,N}=g_{N+1}2^{\,h-1}+\cdots+g_{N+h}$: the
\href{https://github.com/wcook04/plectis-erdos/blob/3d6d938d696fed0fb71dd55115a18a73738ff223/lean/ErdosProblems/Erdos251/PrimeGapDyadicTail.lean\#L647}{tail block}.

\begin{theorem}[block identity]\label{long251:res:block}\compatlabel{res:block}
For every $N$ and $h$,
\[
 T_{N+h}=2^{h}T_N-B_{h,N},
 \qquad\text{hence}\qquad
 \sigma_h(N)=(2^{h}-1)\,T_N-B_{h,N} .
\]
\end{theorem}

\begin{proof}
A routine induction on $h$; the step is one application of the recurrence
together with the recursion defining $B$.  The second identity is the first
minus $T_N$.
\end{proof}

\noindent Formalised as the
\href{https://github.com/wcook04/plectis-erdos/blob/3d6d938d696fed0fb71dd55115a18a73738ff223/lean/ErdosProblems/Erdos251/PrimeGapDyadicTail.lean\#L653}{iterated
block identity} and the
\href{https://github.com/wcook04/plectis-erdos/blob/3d6d938d696fed0fb71dd55115a18a73738ff223/lean/ErdosProblems/Erdos251/PrimeGapDyadicTail.lean\#L667}{scaled
shift identity}.  The shift also obeys the recurrence in its own right,
\begin{equation}\label{long251:eq:shift-step}
 \sigma_h(N+1)=2\sigma_h(N)-\bigl(g_{N+h+1}-g_{N+1}\bigr),
\end{equation}
the
\href{https://github.com/wcook04/plectis-erdos/blob/3d6d938d696fed0fb71dd55115a18a73738ff223/lean/ErdosProblems/Erdos251/PrimeGapDyadicTail.lean\#L563}{shift step
identity}.  Since $B_{h,N}$ is an integer, Theorem~\ref{long251:res:block} converts a
question about the shift into a question about $(2^{h}-1)T_N$: the shift
$\sigma_h(N)$ is integral if and only if $(2^{h}-1)T_N$ is, the
\href{https://github.com/wcook04/plectis-erdos/blob/3d6d938d696fed0fb71dd55115a18a73738ff223/lean/ErdosProblems/Erdos251/PrimeGapDyadicTail.lean\#L802}{integral-shift
criterion}.  The criterion holds in both directions, so once
$T_N$ is fixed no choice of the digits after index $N$ can change whether
$\sigma_h(N)$ is an integer.

For a rational orbit the denominator settles the matter exactly:
\begin{equation}
 \operatorname{den}(T_{N+1})
 =\frac{\operatorname{den}(T_N)}{\gcd(2,\operatorname{den}(T_N))},
 \qquad
 \sigma_h(N)\in\Z
 \iff
 \operatorname{den}(T_N)\mid 2^h-1 .
\label{long251:eq:den-law}\compatlabel{eq:shift-denominator}\compatlabel{res:shiftiff}
\end{equation}
Each even denominator loses exactly one factor of $2$ and an odd denominator is
unchanged, so after finitely many steps the denominator is an odd integer $d$;
Euler's congruence then gives $d\mid2^{\ph(d)}-1$, and the multiplicative order
of $2$ modulo $d$ is the least such exponent.  These are the
\href{https://github.com/wcook04/plectis-erdos/blob/3d6d938d696fed0fb71dd55115a18a73738ff223/lean/ErdosProblems/Erdos251/PrimeGapDyadicTail.lean\#L1217}{denominator
recurrence}, its
\href{https://github.com/wcook04/plectis-erdos/blob/3d6d938d696fed0fb71dd55115a18a73738ff223/lean/ErdosProblems/Erdos251/PrimeGapDyadicTail.lean\#L1226}{odd
case} and
\href{https://github.com/wcook04/plectis-erdos/blob/3d6d938d696fed0fb71dd55115a18a73738ff223/lean/ErdosProblems/Erdos251/PrimeGapDyadicTail.lean\#L1236}{even
case}, the
\href{https://github.com/wcook04/plectis-erdos/blob/3d6d938d696fed0fb71dd55115a18a73738ff223/lean/ErdosProblems/Erdos251/PrimeGapDyadicTail.lean\#L1279}{denominator
classification}, the
\href{https://github.com/wcook04/plectis-erdos/blob/3d6d938d696fed0fb71dd55115a18a73738ff223/lean/ErdosProblems/Erdos251/PrimeGapDyadicTail.lean\#L877}{totient
shift}, and the
\href{https://github.com/wcook04/plectis-erdos/blob/3d6d938d696fed0fb71dd55115a18a73738ff223/lean/ErdosProblems/Erdos251/PrimeGapDyadicTail.lean\#L900}{propagation
of an integral shift to every later index}.  The hypothesis that $d$ is odd
cannot be dropped: if $T_N=1/2$ then $2^{h}-1$ is odd for every $h\ge1$, so no
shift at $N$ is integral.

\begin{theorem}[exact rationality classification]\label{long251:res:escape-irrational}\compatlabel{res:collapse}
Let $T:\N\to\R$ satisfy $T_{N+1}=2T_N-g_{N+1}$ with integer digits $g$.  The
following are equivalent:
\begin{enumerate}[label=\textup{(\roman*)}]
\item $T_0$ is rational;
\item $\sigma_h(N)$ is an integer for some $h\ge1$ and some $N$;
\item for some fixed $h\ge1$, $\sigma_h(N)$ is an integer at every
  sufficiently large $N$.
\end{enumerate}
Consequently $T_0$ is irrational if and only if every positive-length shift is
nonintegral at every index, equivalently if and only if for every $h\ge1$ and
every cutoff some later $N$ has $\sigma_h(N)\notin\Z$.
\end{theorem}

\begin{proof}
For \textup{(i)}$\Rightarrow$\textup{(iii)}, the block identity identifies the
whole orbit with a rational recurrence starting at $T_0$; apply
\eqref{long251:eq:den-law} at an index where the denominator has become odd, and
propagate forward.  The implication
\textup{(iii)}$\Rightarrow$\textup{(ii)} is immediate.  For
\textup{(ii)}$\Rightarrow$\textup{(i)}, the block identity gives
$\sigma_h(N)=(2^h-1)T_N-B_{h,N}$ with $B_{h,N}$ and $\sigma_h(N)$ integers and
$2^h-1\ne0$, so $T_N$ is rational, and $T_N=2^NT_0-B_{N,0}$ then makes $T_0$
rational.  Negating the three conditions gives the two irrationality
formulations.
\end{proof}

\noindent The exact real classifiers are
\href{https://github.com/wcook04/plectis-erdos/blob/3d6d938d696fed0fb71dd55115a18a73738ff223/lean/ErdosProblems/Erdos251/PrimeGapDyadicTail.lean\#L1479}{one
integral positive shift},
\href{https://github.com/wcook04/plectis-erdos/blob/3d6d938d696fed0fb71dd55115a18a73738ff223/lean/ErdosProblems/Erdos251/PrimeGapDyadicTail.lean\#L1525}{eventual
integrality},
\href{https://github.com/wcook04/plectis-erdos/blob/3d6d938d696fed0fb71dd55115a18a73738ff223/lean/ErdosProblems/Erdos251/PrimeGapDyadicTail.lean\#L1551}{pointwise
nonintegrality}, and
\href{https://github.com/wcook04/plectis-erdos/blob/3d6d938d696fed0fb71dd55115a18a73738ff223/lean/ErdosProblems/Erdos251/PrimeGapDyadicTail.lean\#L1572}{cofinal
nonintegrality}.

\begin{lemma}[the boundary condition identifying a true tail]\label{long251:res:true-tail}
Let $a_1,a_2,\ldots$ be real numbers with
$\sum_{j\ge1}|a_j|2^{-j}<\infty$, and let $U_{N+1}=2U_N-a_{N+1}$.
Then
\[
 U_N=\sum_{j\ge1}a_{N+j}2^{-j}\quad\hbox{for every }N
 \quad\Longleftrightarrow\quad 2^{-N}U_N\longrightarrow0.
\]
\end{lemma}
\begin{proof}
Iteration gives
$2^{-N}U_N=U_0-\sum_{j=1}^{N}a_j2^{-j}$.
The limit is zero exactly when $U_0$ equals the full series; subtracting
its first $N$ terms then gives the claimed tail formula.  In particular,
the recurrence alone leaves the homogeneous freedom $C2^N$.
\end{proof}

\paragraph{The actual prime-gap orbit.}
Put
\[
 T_N=\sum_{j\ge1}\frac{g_{N+j}}{2^{\,j}},
\]
which converges by Theorem~\ref{long251:res:infinite}.  The scaled tail equals its
shifted-gap series, the
\href{https://github.com/wcook04/plectis-erdos/blob/3d6d938d696fed0fb71dd55115a18a73738ff223/lean/ErdosProblems/Erdos251/RealPrimeGapTail.lean\#L31}{shifted-gap
series identity}, and satisfies $T_{N+1}=2T_N-g_{N+1}$ with no rationality
hypothesis, the
\href{https://github.com/wcook04/plectis-erdos/blob/3d6d938d696fed0fb71dd55115a18a73738ff223/lean/ErdosProblems/Erdos251/RealPrimeGapTail.lean\#L48}{real
tail recurrence}.  Since $T_0=2S-1=2.349287932\ldots$, the numbers $\Pi$, $S$
and $T_0$ have the same rationality status, and
Theorem~\ref{long251:res:escape-irrational} specialises to the actual series:
irrationality of $\Pi$ is equivalent to cofinal nonintegrality of the positive
shifts of $T$, the
\href{https://github.com/wcook04/plectis-erdos/blob/3d6d938d696fed0fb71dd55115a18a73738ff223/lean/ErdosProblems/Erdos251/RealPrimeGapTail.lean\#L73}{escape
equivalence}.  Nothing about the bridge from the series to its actual tail is
left open; the open content is the escape hypothesis itself.

\subsection{The free-pair criterion and the lcm diagonal}\label{long251:sec:freepair}\compatlabel{sec:diagonal-collapse}

The escape condition quantifies over a fixed shift length.  Freeing the offset
gives an equivalent condition with a weaker-looking producer.

\begin{theorem}[free-pair criterion]\label{long251:res:freepair}
$S$ is irrational if and only if for every $t\ge1$ and every $N_0$ there are
$N,M\ge N_0$ with $M\equiv N\pmod t$ and $T_M-T_N\notin\Z$.
\end{theorem}

\begin{proof}
If $S$ is irrational, take $M=N+t$ for the $N$ supplied by
Theorem~\ref{long251:res:escape-irrational} at shift length $t$.  Conversely, suppose
$S$ is rational.  By \eqref{long251:eq:den-law} the orbit reaches an index $N_0$ beyond
which the reduced denominator is a fixed odd $d$; let $t$ be the
multiplicative order of $2$ modulo $d$.  For $N,M\ge N_0$ with $M\equiv N$
modulo $t$, the difference $T_M-T_N$ is then an integer, contradicting the
stated property at this $t$ and this cutoff.
\end{proof}

\noindent The offset $M-N$ is free, so one nonintegral congruent pair for each
modulus and each cutoff suffices in place of a supply at a single fixed offset.
The mechanism is exact: beyond a rational state with odd reduced denominator
$d$, the difference $T_M-T_N$ is integral precisely when $M\equiv N$ modulo the
order of $2$ in $\Z/d\Z$, the
\href{https://github.com/wcook04/plectis-erdos/blob/3d6d938d696fed0fb71dd55115a18a73738ff223/lean/ErdosProblems/Erdos251/FreePairReduction.lean\#L92}{free-pair
lattice}, so every rational orbit has a cutoff and a modulus at which
integrality holds exactly on the congruence classes, the
\href{https://github.com/wcook04/plectis-erdos/blob/3d6d938d696fed0fb71dd55115a18a73738ff223/lean/ErdosProblems/Erdos251/FreePairReduction.lean\#L105}{lattice
corollary}.  The criterion for the actual series is checked as the
\href{https://github.com/wcook04/plectis-erdos/blob/3d6d938d696fed0fb71dd55115a18a73738ff223/lean/ErdosProblems/Erdos251/FreePairReduction.lean\#L236}{free-pair
equivalence}, over the
\href{https://github.com/wcook04/plectis-erdos/blob/3d6d938d696fed0fb71dd55115a18a73738ff223/lean/ErdosProblems/Erdos251/FreePairReduction.lean\#L186}{actual
real tail recurrence}.  The mechanism is elementary: modulo one the tail
recurrence is multiplication by two, and integrality of $T_M-T_N$ beyond a
rational state is governed by the multiplicative order of $2$ modulo the odd
reduced denominator.

The criterion rules out two shortcuts.  For the actual prime-gap tail,
Section~\ref{long251:sec:xr-compression} combines the prime number theorem with
a first-moment bound to show that, under rationality, some tail value occurs
$\gg_{C,d}X/\log X$ times in $(X,2X]$; here $d$ is the eventual odd reduced
denominator and $C>1$ is fixed.  This is specific to the prime tail, rather
than a consequence of rationality for arbitrary integer-digit orbits, and
equal tails have integral difference, so they do not supply the free-pair
producer.  Also, under the free-pair lattice the difference is an integer
whenever the modulus divides the offset, and an integer cannot lie in
$(\tfrac12,1)$, so the window component of the reduced event below is already
a contradiction on its own.  Both derivations are ordinary proofs, recorded
in the extended record and not formalised.

A second exact reformulation collapses all shift lengths and basepoints onto a
single schedule.  Let $L_0=1$ and $L_j=\lcm(1,\ldots,j)$.

\begin{theorem}[\href{https://github.com/wcook04/plectis-erdos/blob/3d6d938d696fed0fb71dd55115a18a73738ff223/lean/ErdosProblems/Erdos251/OrderLatticeDiagonal.lean\#L153}{lcm-diagonal criterion}]\label{long251:res:lcmdiagonal}\compatlabel{res:lcmdiagonal}
Let $g:\N\to\Z$ and $T:\N\to\R$ satisfy $T_{N+1}=2T_N-g_{N+1}$.  Then
$T_0$ is irrational if and only if
$T_{2L_j}-T_{L_j}\notin\Z$ for every $j\ge0$.
\end{theorem}

\begin{proof}
Theorem~\ref{long251:res:escape-irrational} gives the forward implication.  Conversely,
if $T_0$ is rational then some positive shift length $h$ is integral at every
basepoint beyond an index $N_0$.  Choose $j$ so large that $N_0\le L_j$ and
$h\mid L_j$.  The cocycle identity
$\sigma_{a+b}(N)=\sigma_a(N)+\sigma_b(N+a)$ shows by induction that every
positive multiple of $h$ is integral at every such basepoint, so
$T_{2L_j}-T_{L_j}$ is an integer.
\end{proof}

The factorial schedule has the same divisibility property; the lcm schedule is
pointwise no larger and is the endpoint used here.  Theorems~\ref{long251:res:freepair}
and~\ref{long251:res:lcmdiagonal} are exact reformulations on the class of
integer-digit orbits.  They change the shape of the required producer and they
supply no nonintegrality statement for the actual prime gaps.

\section{A local certificate, and one actual pair}\label{long251:sec:local-certificate}

The classifiers reduce irrationality to a condition of infinite precision
imposed on a complete tail.  Inside $(-1,1)$ integrality is equality with zero,
so on that range the shift step identity \eqref{long251:eq:shift-step} turns
simultaneous integrality of two adjacent shifts into a single comparison of
digits.  What this buys is a certificate attached to one pair of adjacent
indices, whose verification already contradicts integrality there.  The
distance of a shift from the integers never has to be estimated.

\begin{theorem}[adjacent small-shift obstruction]\label{long251:res:smallpair}\compatlabel{res:smallpair}
Let $T$ satisfy the dyadic tail recurrence with integer digits $g$, and fix
$h$ and $N$.  If
\[
 -1<\sigma_h(N)<1,\qquad -1<\sigma_h(N+1)<1,
 \qquad g_{N+h+1}\ne g_{N+1},
\]
then $\sigma_h(N)$ and $\sigma_h(N+1)$ cannot both be integers.  Consequently,
if such a pair occurs beyond every threshold, the $h$-shift is not eventually
integral.
\end{theorem}

\begin{proof}
An integer strictly between $-1$ and $1$ is zero.  If both shifts were
integers, both would vanish, and \eqref{long251:eq:shift-step} would give
$g_{N+h+1}=g_{N+1}$.  The cofinal statement chooses one such adjacent pair
after the alleged onset of integrality.
\end{proof}

\begin{corollary}[real form and the sufficient condition]\label{long251:res:smallpair-real}\compatlabel{res:smallpair-real}
The same statement holds for a real orbit, with the same proof.  If for every
$h\ge1$ and every cutoff some later $N$ satisfies the three displayed
conditions for the actual prime gaps, then $\Pi$ is irrational.
\end{corollary}

\noindent The rational finite contradiction is the
\href{https://github.com/wcook04/plectis-erdos/blob/3d6d938d696fed0fb71dd55115a18a73738ff223/lean/ErdosProblems/Erdos251/PrimeGapDyadicTail.lean\#L979}{adjacent
small-shift obstruction} with its
\href{https://github.com/wcook04/plectis-erdos/blob/3d6d938d696fed0fb71dd55115a18a73738ff223/lean/ErdosProblems/Erdos251/PrimeGapDyadicTail.lean\#L1006}{cofinal
form}; the real form is the
\href{https://github.com/wcook04/plectis-erdos/blob/3d6d938d696fed0fb71dd55115a18a73738ff223/lean/ErdosProblems/Erdos251/RealPrimeGapTail.lean\#L99}{real
adjacent-pair consumer}, and the implication to irrationality of the actual
series is the
\href{https://github.com/wcook04/plectis-erdos/blob/3d6d938d696fed0fb71dd55115a18a73738ff223/lean/ErdosProblems/Erdos251/RealPrimeGapTail.lean\#L120}{small-mismatch
criterion}.  The third hypothesis alone is available for the actual gaps: it
reads $g_4=2\ne4=g_3$ at $h=1$, $N=2$, and by
Proposition~\ref{long251:res:gap-nonperiodic} it holds for arbitrarily large $N$ at
every fixed $h$.  What is missing is the pair of inequalities, each of which
constrains a complete infinite tail.

\begin{proposition}[prime gaps do not become periodic]\label{long251:res:gap-nonperiodic}\label{res:gap-nonperiodic}
For every positive $h$, the actual consecutive-prime-gap sequence is not
eventually periodic with period $h$.
\end{proposition}

\begin{proof}
For $m\ge2$ the integers $(m+1)!+2,\ldots,(m+1)!+m+1$ are composite, so the
gaps are unbounded.  An eventually periodic sequence of natural numbers has
finite range after its preperiod and a bounded initial segment, hence is
bounded.
\end{proof}

\noindent Lean checks the factorial argument as
\href{https://github.com/wcook04/plectis-erdos/blob/3d6d938d696fed0fb71dd55115a18a73738ff223/lean/ErdosProblems/Erdos251/PrimeGapDyadicTail.lean\#L57}{unboundedness
of the actual gaps} and the conclusion as
\href{https://github.com/wcook04/plectis-erdos/blob/3d6d938d696fed0fb71dd55115a18a73738ff223/lean/ErdosProblems/Erdos251/PrimeGapDyadicTail.lean\#L1023}{non-eventual
periodicity}.  Far stronger lower bounds for large gaps are
known~\cite[Theorem~1]{fgkmt2018}; unboundedness is all that is needed.

The conjunction in Theorem~\ref{long251:res:smallpair} has an exact normal form.  For
fixed $h$ write $D_N=\sigma_h(N)$ and $\delta_N=g_{N+h+1}-g_{N+1}$, so that
$D_{N+1}=2D_N-\delta_N$.

\begin{theorem}[\href{https://github.com/wcook04/plectis-erdos/blob/3d6d938d696fed0fb71dd55115a18a73738ff223/lean/ErdosProblems/Erdos251/AffineShiftEscape.lean\#L113}{signed two-window normal form}]\label{long251:res:signedwindow}
Assume $\delta_N$ is even.  The conjunction
$-1<D_N<1$, $-1<D_{N+1}<1$, $\delta_N\ne0$ is equivalent to
\[
 \bigl(\delta_N=2\ \hbox{ and }\tfrac12<D_N<1\bigr)
 \quad\hbox{or}\quad
 \bigl(\delta_N=-2\ \hbox{ and }-1<D_N<-\tfrac12\bigr).
\]
\end{theorem}

\begin{proof}
The recurrence and the two unit windows give $-3<\delta_N<3$.  A nonzero even
integer in that interval is $2$ or $-2$.  Substitution into
$D_{N+1}=2D_N-\delta_N$ gives the stated half-window and, in the reverse
direction, recovers the second unit window.
\end{proof}

For the actual gaps and $N\ge1$ every $\delta_N$ is even, so
Theorem~\ref{long251:res:signedwindow} applies throughout the range where the local
certificate is sought.  The same computation for a real orbit is one line and
is used in Theorem~\ref{long251:res:sparse} below.

\subsection{An explicit remainder and a certified actual pair}

Both window conditions involve complete infinite tails.  A dominated
truncation reduces each of them to a finite integer comparison.

\begin{proposition}[explicit remainder]\label{long251:res:explicit-remainder}
For integers $h,N\ge0$ and $L\ge1$ put
\[
 F_{h,N,L}=\sum_{j=1}^{L}\frac{g_{N+h+j}-g_{N+j}}{2^{\,j}},\qquad
 P(x)=x^4+8x^3+36x^2+104x+150,
\]
\[
 E_{h,N,L}=\frac{1250}{2^{L}}\bigl(P(N+h+L+2)+P(N+L+2)\bigr).
\]
Then $\bigl|\sigma_h(N)-F_{h,N,L}\bigr|\le E_{h,N,L}$.  In particular
$|F_{h,N,L}|+E_{h,N,L}<1$ certifies $|\sigma_h(N)|<1$, and
$\operatorname{dist}(F_{h,N,L},\Z)>E_{h,N,L}$ certifies
$\sigma_h(N)\notin\Z$.
\end{proposition}

\begin{proof}
The checked bound $p_n\le1250(n+1)^4$ gives $0\le g_m\le p_{m+1}\le1250(m+2)^4$,
so the omitted absolute tail is at most
\[
 1250\sum_{j>L}\frac{(N+h+j+2)^4+(N+j+2)^4}{2^{\,j}} .
\]
The polynomial identity $2P(x)=(x+1)^4+P(x+1)$ telescopes to
$\sum_{k=1}^{J}(m+k)^42^{-k}=P(m)-P(m+J)2^{-J}$, whose last term tends to zero.
Apply it with $m=N+h+L+2$ and $m=N+L+2$ after writing $j=L+k$.  The two tests
follow from the triangle inequality and the definition of distance to $\Z$.
\end{proof}

\begin{proposition}[a certified adjacent pair]\label{long251:res:finite-smallpair}
For the actual prime gaps, $h=1$ and $N=2$ satisfy the three hypotheses of
Theorem~\ref{long251:res:smallpair}: both $\sigma_1(2)$ and $\sigma_1(3)$ lie in
$(-1,1)$ and are nonintegral, and $g_4=2\ne4=g_3$.
\end{proposition}

\begin{proof}
Take $L=40$ and $Q=2^{40}=1099511627776$.  Exact integer arithmetic over the
first $46$ primes gives
\[
\begin{array}{c|r|r}
 N&Q\,F_{1,N,40}&Q\,E_{1,N,40}\\\hline
 2&-662838684750&11764181250\\
 3& 873345886050&12805761250
\end{array}
\]
In each row $|QF|+QE<Q$, which puts the whole certified interval inside
$(-1,1)$, and $QE$ is smaller than the distance from $QF$ to the nearest
multiple of $Q$, which excludes every integer.  The margins are $424908761776$
and $213359980476$ respectively.  These are strict integer comparisons rather
than inferences from rounded numerical tails.
Appendix~\ref{long251:app:certificates} replays them.
\end{proof}

This witness proves the local condition at one pair of indices.  The
quantifiers over every $h$ and every cutoff remain.

\begin{proposition}[a one-tail signed certificate]\label{long251:res:one-tail-certificate}
Let $D_{N+1}=2D_N-\delta_N$ with real $D_N$, and suppose
$\delta_N=2s$ for $s\in\{-1,1\}$.  If integers $A,B,Q$ satisfy
$Q>0$, $B\ge0$, $|QD_N-A|\le B$, and
\[
 2sA-Q>2B,\qquad Q-sA>B,
\]
then $1/2<sD_N<1$, $|D_{N+1}|<1$, and both $D_N,D_{N+1}$ are
nonintegral.  In particular the adjacent small-shift obstruction is
certified from just one tail enclosure.
\end{proposition}
\begin{proof}
The inequalities place the complete interval
$[(sA-B)/Q,(sA+B)/Q]$ inside $(1/2,1)$.
Thus $sD_{N+1}=2sD_N-2\in(-1,0)$, proving every assertion.
\end{proof}
For $h=1,N=2$, the exact data already used above give
$Q=2^{40}$, $A=-662838684750$, $B=11764181250$ and $s=-1$.
The two positive integer margins are respectively
$202637379224$ and $424908761776$.
The second enclosure in Proposition~\ref{long251:res:finite-smallpair}
remains a useful independent cross-check, not a logical necessity.

\section{Two denominator floors}\label{long251:sec:denominator-floor}

The same polynomial bound yields an exact rational enclosure of $\Pi$, and any
such enclosure excludes an initial range of denominators.

\begin{theorem}[kernel-decided denominator floor]\label{long251:res:denominatorfloor}
Let $a\in\Z$ and $b\in\Npos$.  If $\Pi=a/b$ then $b\ge2^{589}>10^{177}$, and
the same floor holds for every rational equal to $S$.
\end{theorem}

\begin{proof}
Put $c=1229$, the number of primes below $10^4$, and
$A=\sum_{i=0}^{c-1}p_i2^{\,c-i-1}$, so that $A/2^{c}$ is the $c$th partial sum
of $\Pi$.  All terms are positive, so $A/2^{c}\le\Pi$.  For the upper bound,
$p_{c+j}\le1250(c+j+1)^4$ and $(c+1+j)^4\le(c+1)^4(3/2)^{j}$ for $j\ge0$, valid
because $c\ge9$; summing the resulting geometric tail with ratio $3/4$ gives
\[
 \frac{A}{2^{c}}\;\le\;\Pi\;\le\;\frac{2A+5000(c+1)^4}{2^{\,c+1}} .
\]
The four positive integers $u,v,u',v'$ printed in
Appendix~\ref{long251:app:certificates} satisfy
\[
 u'v-uv'=1,\qquad
 \frac uv<\frac{A}{2^{c}},\qquad
 \frac{2A+5000(c+1)^4}{2^{\,c+1}}<\frac{u'}{v'},\qquad
 v+v'\ge2^{589},
\]
all four being integer comparisons after clearing the positive denominators.
If $\Pi=a/b$ then $av-bu\ge1$ and $bu'-av'\ge1$, and the determinant identity
gives
\[
 b=b(u'v-uv')=(av-bu)v'+(bu'-av')v\ge v+v' \ge 2^{589}.
\]
This is the classical bound for a rational strictly between two Farey
neighbours.
Since $589\log_{10}2>177$, the floor exceeds $10^{177}$.  If $S=a/b$, apply the
same argument to $\Pi=(a+2b)/b$.
\end{proof}

\noindent The Lean kernel decides the four inequalities on the same literals by
\texttt{decide +kernel} with no \texttt{native\_decide}, over a trial-division
sieve it re-runs on $[2,10^4)$: the
\href{https://github.com/wcook04/plectis-erdos/blob/3d6d938d696fed0fb71dd55115a18a73738ff223/lean/ErdosProblems/Erdos251/KernelDenominatorFloor.lean\#L310}{certificate},
the
\href{https://github.com/wcook04/plectis-erdos/blob/3d6d938d696fed0fb71dd55115a18a73738ff223/lean/ErdosProblems/Erdos251/KernelDenominatorFloor.lean\#L316}{floor
for the prime series}, and the
\href{https://github.com/wcook04/plectis-erdos/blob/3d6d938d696fed0fb71dd55115a18a73738ff223/lean/ErdosProblems/Erdos251/KernelDenominatorFloor.lean\#L324}{floor
for the gap series}.  Its analytic input is the same enclosure, over the
polynomial bound at
\href{https://github.com/wcook04/plectis-erdos/blob/3d6d938d696fed0fb71dd55115a18a73738ff223/lean/ErdosProblems/Erdos251/PrimeGapDyadicTail.lean\#L360}{line~360}.
The four Farey literals have $177$ and $178$ digits.

\begin{theorem}[certified continued-fraction exclusion]\label{long251:res:cfexclusion}
Every rational equal to $\Pi$, and hence every rational equal to $S$, has
reduced denominator $q\ge2^{39997}$, and therefore $q>10^{12040}$.
\end{theorem}

\emph{Evidence.}  This is an exact finite computation and no Lean declaration
carries it.  The continued-fraction expansion of $\Pi$ is developed to $23369$
partial quotients at a working scale of $80000$ bits.  Each quotient is forced
by a rational bracket of width $1$ in those scaled integer units around the
true value, and the separation $|q_n\Pi-p_n|>0$ is verified against that
bracket, so the prefix is the true prefix.  No floating-point reconstruction
enters it.  The bracket uses the Lean-checked tail estimate
$p_n\le1250(n+1)^4$, so no conjecture enters.  The classical theorem that a
best approximation of the second kind is a convergent then converts the
certified prefix into the denominator bound
\cite[\S6, Theorems~16 and~17]{khinchin1964}; Short gives an accessible
statement and proof~\cite[Theorem~1.1]{short2009}.  The receipt
is \href{https://github.com/wcook04/plectis-erdos/blob/dd6cc708f650f061de06594eba6ba43878f4de31/research/experiments/erdos251/receipts/certified-cf.json}{\texttt{certified-cf.json}},
which records the generating program and its digest, the power-of-two exponent
$39997$, the bit length $39998$, the strict decimal power $12040$, the decimal
digit count $12041$, the largest partial quotient $973919$ at index $16442$,
and an arithmetic self-check of the same routine against the known expansions
of $e$ and $\pi$.  The certified statement is $q\ge2^{39997}$ and
$q>10^{12040}$.  Printing the bit length or the digit count as an exponent
overstates the bound by one power of two and one decimal order.

Each floor is finite, and extending either one raises the excluded range and
reproduces the same statement form.  Neither decides the problem.

\section{What cannot supply the missing input}\label{long251:sec:obstructions}

The counterexamples in this section test which information about the gaps
could force irrationality.  They preserve selected growth, residue and
nonconcentration properties while changing the dyadic value.  The sparsity
result has a different role: it rules out a positive-proportion lower bound
for the two-window event, while leaving sparse-witness arguments available.

\subsection{Bounded residue-preserving perturbations}

\begin{theorem}[bounded-perturbation obstruction]\label{long251:res:boundedperturbation}
Let $a_n$ be natural numbers with $\sum_{n\ge0}a_n2^{-(n+1)}$ convergent.  For
every integer $M\ge1$ and every cutoff $K$ there are digits
$\varepsilon_n\in\{0,1\}$, zero for $n<K$, such that
$\sum_{n\ge0}(a_n+M\varepsilon_n)2^{-(n+1)}$ is rational.
\end{theorem}

\begin{proof}
Write $A=\sum_{n\ge0}a_n2^{-(n+1)}$ and choose a rational
$r\in\bigl(A,A+M2^{-K}\bigr)$.  A binary expansion of $(r-A)/M$ has the form
$\sum_{n\ge K}\varepsilon_n2^{-(n+1)}$ with $\varepsilon_n\in\{0,1\}$; set
$\varepsilon_n=0$ for $n<K$.  The perturbed series converges and equals $r$.
\end{proof}

\noindent The
\href{https://github.com/wcook04/plectis-erdos/blob/3d6d938d696fed0fb71dd55115a18a73738ff223/lean/ErdosProblems/Erdos251/BoundedPerturbationCountermodel.lean\#L193}{construction}
is kernel checked.  The invariant-property consequence---a property shared by
every admissible perturbation cannot by itself force irrationality---is the
ordinary logical reading of that construction, rather than a separately named
Lean theorem.  The construction is also
checked at the actual consecutive prime gaps for every $M\ge1$ and every $K$,
the
\href{https://github.com/wcook04/plectis-erdos/blob/3d6d938d696fed0fb71dd55115a18a73738ff223/lean/ErdosProblems/Erdos251/BoundedPerturbationCountermodel.lean\#L208}{rational
perturbed prime-gap series}, with the perturbed gap lying in $[g_n,g_n+M]$ and
congruent to $g_n$ modulo $M$, the
\href{https://github.com/wcook04/plectis-erdos/blob/3d6d938d696fed0fb71dd55115a18a73738ff223/lean/ErdosProblems/Erdos251/BoundedPerturbationCountermodel.lean\#L234}{gap
bounds}.  Consequently no theorem about the size or the residue of prime gaps
that survives such a perturbation can suffice for Problem~\#251; that class
includes the bounded-gap and large-gap theorems cited in
Section~\ref{long251:sec:open} and equidistribution of $p_n$ modulo any fixed $q$, on
taking $M=2q$.  The perturbed digits need not remain prime gaps.

\subsection{Algebraic nonconcentration survives the rationalising perturbation}

Fixed-block polynomial nonconcentration is a finer property than size and
residue, but it too survives the bounded perturbation.  The relevant input
is Schlage-Puchta's Lemma~4: for every polynomial
$F\in\Z[x_0,\ldots,x_k]$ which does not vanish identically,
$F(g_n,\ldots,g_{n+k})\ne0$ for almost all $n$~\cite[Lemma~4,
pp.~5--6]{schlagepuchta2011}.  Say that an integer sequence $a$ has
\emph{fixed-block nonconcentration} when it satisfies that conclusion for
every $k\ge0$ and every nonzero $F\in\Z[x_0,\ldots,x_k]$.  The property
transfers across bounded perturbations with no independence or randomness
assumption.


\begin{proposition}[finite counting for shifted gap differences]\label{long251:res:shiftedcount}
Write $p_n$ for the primes indexed from $p_0=2$ and $g_n=p_{n+1}-p_n$.
For $h\ge2$ and $r\in\Z$, let $M_{h,r}(N)$ count $n<N$ with $g_{n+h}-g_n=r$.
Let $Q_{N,H,r}$ count triples $(x,d,s)$ with $x<p_N$,
$0<d<s\le H$, $d+r>0$, and all four integers
$x,x+d,x+s,x+s+d+r$ prime. Then, for every $N,H\ge0$,
\[
 (H+1)M_{h,r}(N)\le
 (h+1)p_{N+h+1}+(H+1)Q_{N,H,r}.
\]
\end{proposition}

\begin{proof}
For a counted index $n$, put $x=p_n$, $d=g_n$ and
$s=p_{n+h}-p_n$.  If the total span $p_{n+h+1}-p_n$ is at most $H$,
then $0<d<s\le H$, $d+r=g_{n+h}>0$, and the four primes are precisely
$x,x+d,x+s,x+s+d+r$.  The map is injective because $x=p_n$ determines
$n$.  For the remaining indices the integer span is at least $H+1$.
Each gap appears in at most $h+1$ of the spans, so their total is at most
$(h+1)\sum_{i<N+h}g_i<(h+1)p_{N+h+1}$.  Multiply the resulting bound
on the exceptional count by $H+1$ and add the small-span contribution.
\end{proof}

The displayed finite shifted-count bound is unconditional by the ordinary
proof above.  Its \href{https://github.com/wcook04/plectis-erdos/blob/d058b9150218ae23d2cfd975580088d6167e9da8/ErdosProblems/Erdos251/ShiftedGapCountingR9.lean#L152}{historical source locator}
is not present in the supplied declaration index, so no checked-build status
is assigned to that link here.
To deduce zero density one still needs, for every $\varepsilon>0$ and all
large $N$, a choice of $H$ making the right side less than
$\varepsilon N(H+1)$. The adjacent case $h=1$ instead requires the
corresponding three-prime estimate, because two of the four positions would
coincide. The all-shift implication records these analytic premises explicitly;
it supplies neither estimate and does not prove irrationality.

\begin{theorem}[nonconcentration is perturbation-stable]\label{long251:res:nonconcentration}
Let $a:\N\to\Z$ have fixed-block nonconcentration, let $E\subset\Z$ be finite,
and let $b_n=a_n+e_n$ with $e_n\in E$ for every $n$.  Then $b$ has fixed-block
nonconcentration.
\end{theorem}

\begin{proof}
Fix $k\ge0$ and a nonzero $F\in\Z[x_0,\ldots,x_k]$.  For a tuple
$\mathbf e=(e^{(0)},\ldots,e^{(k)})\in E^{k+1}$ put
$F_{\mathbf e}(x_0,\ldots,x_k)=F(x_0+e^{(0)},\ldots,x_k+e^{(k)})$.  The
substitution $x_i\mapsto x_i+e^{(i)}$ is a ring automorphism of
$\Z[x_0,\ldots,x_k]$ with inverse $x_i\mapsto x_i-e^{(i)}$, so
$F_{\mathbf e}\ne0$.  If $F(b_n,\ldots,b_{n+k})=0$ then
$F_{\mathbf e}(a_n,\ldots,a_{n+k})=0$ for the tuple
$\mathbf e=(e_n,\ldots,e_{n+k})$, whence
\[
 \{\,n:F(b_n,\ldots,b_{n+k})=0\,\}
 \;\subseteq\!\!
 \bigcup_{\mathbf e\in E^{k+1}}\!\!
 \{\,n:F_{\mathbf e}(a_n,\ldots,a_{n+k})=0\,\}.
\]
Each set on the right has density zero by hypothesis, and the union is over
the finite index set $E^{k+1}$, so a finite union of density-zero sets has
density zero.
\end{proof}

\begin{corollary}[nonconcentration does not force irrationality]\label{long251:res:nonconc-primes}
Fix $M\ge1$ and $K\ge0$, and let $b$ be the perturbed sequence supplied by
Theorem~\ref{long251:res:boundedperturbation} at the actual prime gaps.  Then
$\sum_{n\ge0}b_n2^{-(n+1)}$ is rational, $b_n=g_n$ for $n<K$,
$b_n-g_n\in\{0,M\}$ and $b_n\equiv g_n\pmod M$ for every $n$, $b$ has
fixed-block nonconcentration, and the cumulative sequence
$P_n=2+\sum_{i<n}b_i$ satisfies $p_n\le P_n\le p_n+Mn$ and hence
$P_n\sim n\log n$.
\end{corollary}

\begin{proof}
Everything except nonconcentration is Theorem~\ref{long251:res:boundedperturbation}
together with $\sum_{i<n}g_i=p_n-2$, and $P_n\sim n\log n$ follows from
$p_n\sim n\log n$~\cite[Chapter~6]{mv2007}.  The perturbation takes values in
the fixed two-element set $E=\{0,M\}$, so
Theorem~\ref{long251:res:nonconcentration} applies to the actual gaps, which
have fixed-block nonconcentration by Schlage-Puchta's
lemma~\cite[Lemma~4]{schlagepuchta2011}.
\end{proof}

Fixed-block polynomial nonconcentration, taken together with a prescribed
finite prefix of actual gaps, a pointwise bound $b_n\le g_n+M$, every residue
modulo $M$, positivity and cumulative growth at the prime-number-theorem scale,
is therefore compatible with a rational dyadic value.  Taking $M$ even also
preserves eventual evenness.  To preserve a prescribed modulus $q$ together
with parity, take $M=2q$, with the corresponding pointwise bound
$b_n\le g_n+2q$.  The boundary is exact: the terms $P_n$ are not asserted to be prime,
and the conclusion says nothing about constraints with a block length or a
polynomial family growing with the index.

\begin{proposition}[nonconcentration under sparse changes]\label{long251:res:sparse-nonconcentration}
Let $a,b:\N\to\Z$ agree off a set $S$ of ordinary density zero.
If $a$ has fixed-block nonconcentration, then so does $b$.
No boundedness assumption on $a-b$ is needed.
\end{proposition}
\begin{proof}
For fixed $k$ and a nonzero $F\in\Z[x_0,\ldots,x_k]$, a zero of
$F(b_n,\ldots,b_{n+k})$ either is already a zero at the $a$-block or has
$n+i\in S$ for some $0\le i\le k$.  Each fixed translate of a
density-zero set has density zero, and a finite union retains that
property.  This proves the assertion for each fixed $F$ and $k$.
It gives no uniform estimate for polynomial families or lengths growing
with $n$.
\end{proof}

\begin{corollary}[simultaneous prime-gap countermodel]\label{long251:res:jointcountermodel}
For every prescribed finite prime-gap prefix and $0<\varepsilon\le1$,
there is an altered word $b=g+e$, with $P_n=2+\sum_{i<n}b_i$,
for which one can simultaneously impose a rational dyadic value, nonnegative
integer corrections eventually at most $(\log(n+3))^\varepsilon$,
all fixed eventual coefficient and cumulative congruences,
fixed-block polynomial nonconcentration, and vanishing literal block
TV distance for lengths $o(\log\log X)$.  The cumulative positions satisfy
\[
 0\le P_n-p_n=O_\varepsilon\left(
    \frac{n(\log(n+3))^\varepsilon}{\log\log n}\right),
 \qquad P_n\sim n\log n.
\]
\end{corollary}
\begin{proof}
Apply the sparse theorem at a rational target, then the preceding
proposition and Schlage-Puchta's Lemma~4 \cite{schlagepuchta2011}.
To obtain the cumulative support bound, split $[\sqrt n,n)$ into dyadic
intervals, use $\log\log x\asymp\log\log n$ there, and bound the first
$\sqrt n$ indices trivially.  Thus $|S\cap[0,n)|=O(n/\log\log n)$.
The pointwise correction bound gives the displayed estimate, including
a fixed constant for the exceptional prefix.  Since
$(\log n)^{\varepsilon-1}/\log\log n\to0$ for $0<\varepsilon\le1$,
the prime number theorem \cite[Chapter~6]{mv2007} gives the last conclusion.
\end{proof}
These are ordinary deductions from the sparse theorem and the stated
external analytic input.  The reconstructed positions need not be prime.
Neither a bounded translation nor a sparse change preserves general
relative frequencies of rare events; their two transfer arguments have
different hypotheses and must not be conflated.

\subsection{The two-window event has density zero}

The same lemma bounds the event that the local certificate consumes.

\begin{theorem}[sparsity of the two-window event]\label{long251:res:sparse}
Fix $h\ge1$.  The set of $N\ge1$ at which the three hypotheses of
Theorem~\ref{long251:res:smallpair} hold for the actual prime gaps has density zero.
For the same $h$, the set of $N$ with $g_{N+h+1}=g_{N+1}$ also has density
zero.
\end{theorem}

\begin{proof}
For $N\ge1$ every gap involved is even, so $\delta_N=g_{N+h+1}-g_{N+1}$ is
even.  If $|\sigma_h(N)|<1$, $|\sigma_h(N+1)|<1$ and $\delta_N\ne0$, then
$\delta_N=2\sigma_h(N)-\sigma_h(N+1)$ lies in $(-3,3)$, so $\delta_N=\pm2$.
Apply Schlage-Puchta's lemma~\cite[Lemma~4]{schlagepuchta2011} at index
$n=N+1$ with $k=h$ to the two
polynomials $F_\pm(x_0,\ldots,x_h)=x_h-x_0\mp2$, neither of which vanishes
identically: each of $\{N:\delta_N=2\}$ and $\{N:\delta_N=-2\}$ has density
zero, and the event is contained in their union.  The second assertion is the
same lemma applied to $F(x_0,\ldots,x_h)=x_h-x_0$.
\end{proof}

The two halves point in opposite directions and both are useful.  The mismatch
hypothesis on its own holds for almost all $N$, which is strictly stronger than
the cofinal statement Proposition~\ref{long251:res:gap-nonperiodic} supplies.  The full
conjunction is confined to a set of density zero, because the two window
conditions force the mismatch to be exactly $\pm2$.  A producer for
Problem~\ref{long251:prob:smallpair} therefore has to be a cofinality statement about
a sparse set.  Consequently, this event cannot be supplied on a set of positive
lower density.  A sufficient producer would have to yield cofinally many
witnesses in a density-zero set; density zero alone does not exclude an
averaging argument capable of detecting such sparse witnesses.  The finite measurements in Section~\ref{long251:sec:measurements} show a
decline across the sampled bands.  They establish neither an asymptotic
rate for the density nor a counting law.  Prime-height cutoffs and gap-index
cutoffs must also be distinguished before comparing any proposed rates.

\subsection{Recurring gap values differing by two do not suffice}

Theorem~\ref{long251:res:signedwindow} reduces the producer to a $\pm2$ mismatch
together with a half-window condition.  It is tempting to hope that the
mismatch half is the substance, and that two even values differing by $2$, each
occurring infinitely often as a consecutive prime gap inside a common index
residue class, would suffice.  At the level of integer-digit recurrences that
implication is false, and it stays false under growth at the scale the primes
actually have.

\begin{theorem}[recurring values are not enough]\label{long251:res:polignacfail}
There is a sequence $(a_n)_{n\ge1}$ of positive even integers with the
following properties.  The values $2$ and $4$ each occur infinitely often at
indices divisible by every fixed $t\ge1$.  The sequence is unbounded, not
eventually periodic, and satisfies $a_n=O(\log n)$.  The series
$\sum_{n\ge1}a_n2^{-n}$ equals $6$, and every scaled tail
$\sum_{j\ge1}a_{N+j}2^{-j}$ is an integer, so every tail shift is integral.
The increasing odd sequence $P_n=3+\sum_{j\le n}a_j$ satisfies
$P_n\sim n\log n$.
\end{theorem}

\begin{proof}
Let logarithms be natural and put $b_n=2\lceil\tfrac12\log(n+64)\rceil$, so
that $b_n$ is even, $b_n\ge6$, and $b_n=\log(n+64)+O(1)$.  Since
$\tfrac12\log(n+65)-\tfrac12\log(n+63)<1$ for every $n\ge1$, consecutive
increments of $b$ lie in $\{0,2\}$ and $b_{n+1}-b_{n-1}\le2$.  Call an index
$n$ \emph{special} when $n=k!$ or $n=2\,k!$ for some $k\ge5$, and define
\[
 U_n=\begin{cases}
  2b_{n-1}-2,& n=k!,\ k\ge5,\\
  2b_{n-1}-4,& n=2\,k!,\ k\ge5,\\
  b_n,&\text{otherwise},
 \end{cases}
 \qquad a_n=2U_{n-1}-U_n\quad(n\ge1).
\]
The special indices are distinct and never adjacent, since $k!$ and $2\,k!$ are
even and differ by more than one from each other and from the neighbouring
special indices.  Each $U_n$ is an even integer by construction, so
each $a_n$ is an even integer.

At an ordinary index whose predecessor is also ordinary,
$a_n=2b_{n-1}-b_n=b_{n-1}-(b_n-b_{n-1})\ge6-2=4$.  At a special index the
predecessor is ordinary and $a_n=2b_{n-1}-(2b_{n-1}-r)=r$, which is $2$ at
$n=k!$ and $4$ at $n=2\,k!$.  Immediately after a special index carrying $r$,
\[
 a_{n+1}=2(2b_{n-1}-r)-b_{n+1}\ge4b_{n-1}-2r-(b_{n-1}+2)
 =3b_{n-1}-2r-2\ge8 .
\]
Every $a_n$ is therefore a positive even integer, and $U_n=O(\log n)$ gives
$a_n=O(\log n)$.  At ordinary indices with ordinary predecessors
$a_n\ge b_{n-1}-2\to\infty$, so $(a_n)$ is unbounded and hence not eventually
periodic.

The definition $a_n=2U_{n-1}-U_n$ is exactly $a_n2^{-n}=U_{n-1}2^{-(n-1)}-U_n2^{-n}$,
so for every $N\ge0$ and $m\ge1$
\[
 \sum_{j=1}^{m}\frac{a_{N+j}}{2^{\,j}}=U_N-\frac{U_{N+m}}{2^{\,m}} .
\]
Since $U_n=O(\log n)$ the endpoint tends to zero, so the tail at $N$ equals the
integer $U_N$; at $N=0$ the series equals $U_0=b_0=6$.  Every difference
$U_{N+h}-U_N$ is an integer, so every tail shift is integral.

For each $t$ and every sufficiently large $k$ we have $t\mid k!$, hence
$t\mid2\,k!$, and $a_{k!}=2$ while $a_{2k!}=4$: both values recur infinitely
often at indices congruent to $0$ modulo $t$.

Finally, summing $a_j=2U_{j-1}-U_j$ gives
$\sum_{j=1}^{n}a_j=2U_0+\sum_{j=1}^{n-1}U_j-U_n$.  The baseline
$\sum_{j<n}b_j=n\log n+O(n)$.  The special indices up to $n$ number
$O(\log n/\log\log n)$ and each changes the summand by $O(\log n)$, so their
total contribution is $O((\log n)^2)=o(n)$.  Hence
$P_n=3+\sum_{j\le n}a_j\sim n\log n$, and $P_n$ is odd and increasing.
\end{proof}

The factorial schedule above proves recurrence in the zero residue class.
The following separate construction strengthens this to every residue class.


\subsection{A complete countermodel in every residue class}
\label{long251:res:all-residue-log-countermodel}
There is a synthetic integer sequence $(a_n)_{n\ge1}$ with all of the
following properties.  Each $a_n$ is positive and divisible by $2$, and
\[
 a_n\le 4\log(n+1)+24.
\]
For every $t\ge1$, every $0\le r<t$, and every cutoff $N$, there are
$i,j\ge N$ such that
\[
 i\equiv j\equiv r\pmod t,\qquad a_i=2,\quad a_j=4.
\]
For every integer $B$ and cutoff $N$, some $n\ge N$ satisfies $a_n>B$.
For every $h\ge1$ there is no cutoff after which $a_{n+h}=a_n$ for all $n$.
Nevertheless, all the complete tails
\[
 U_N=\sum_{j\ge1}\frac{a_{N+j}}{2^j}\quad(N\ge0)
\]
are integers, $U_0=6$, and $U_{N+h}-U_N\in\mathbb Z$ for every $N,h\ge0$.
The sequence
\[
 P_N=3+\sum_{j=1}^{N}a_j
\]
is strictly increasing and odd, with $P_N/(N\log N)\to1$.
These are synthetic positions; no $P_N$ is asserted to be prime.

To obtain every residue class, enumerate triples consisting of a positive
modulus, a residue, and a repetition index.  Choose the corresponding
centres $c_k$ in the prescribed classes with $c_0\ge100$,
$c_{k+1}\ge c_k+3$ and $c_k\ge2^{k^2}$ for $k\ge1$.
The parity of the repetition index selects the desired value $v_k=2$ or $4$.
Use the even baseline
\[
 b_n=6+2\left\lfloor\frac{\log(n+1)}2\right\rfloor,
 \qquad
 U_n=\begin{cases}2b_{n-1}-v_k,&n=c_k,\\b_n,&n\notin\{c_k:k\ge0\},\end{cases}
\]
and define $a_n=2U_{n-1}-U_n$.  The separation of the centres and the
baseline increment bound give positivity and the displayed logarithmic
bound.  Finite telescoping leaves $U_{N+m}/2^m$, which tends to zero;
thus these carries are the complete tails.  The sparse changes to the
baseline have negligible contribution after division by $N\log N$,
which gives the stated growth of $P_N$.

The full statement is \href{https://github.com/wcook04/plectis-erdos/blob/3d6d938d696fed0fb71dd55115a18a73738ff223/lean/ErdosProblems/Erdos251/AllResidueLogarithmicR9.lean\#L503}{kernel-checked in Lean}.


This rules out an inference from these coefficient properties alone to
nonintegral tail shifts.  It gives no counterexample to the prime-gap problem
and does not reproduce the extreme large gaps of the primes.

The consequence is exact.  Positivity, evenness, unboundedness, non-eventual
periodicity, a pointwise logarithmic bound, cumulative growth at the
prime-number-theorem scale, and the recurrence of two even values differing by
$2$ inside every index residue class are jointly compatible with an integral
tail at every index.  Any route through that input must use a property of the
consecutive primes beyond this list.  The sequence $(a_n)$ is synthetic and the
$P_n$ are not asserted to be prime.  The pointwise logarithmic bound also
precludes the known large-gap behaviour, so the theorem does not speak to
hypotheses that use extreme gaps.

\subsection{Two further boundaries}

\begin{proposition}[quadratic polynomial-shift countermodel]\label{long251:res:polynomialcountermodel}\compatlabel{res:polynomialcountermodel}
Put $c_n=2(n^2+4n+2)$ and $U_n=2(n+4)^2$.  Then $c_n$ is positive, even and
strictly increasing, $U_{n+1}=2U_n-c_{n+1}$, every shift $U_{N+h}-U_N$ is
integral, $c_{n+1}-c_n=4n+10$ is never $\pm2$, and
\[
 \sum_{j\ge1}\frac{c_j}{2^{\,j}}=32 .
\]
\end{proposition}

\begin{proof}
Direct expansion gives the recurrence, and the finite telescope is
$\sum_{j=1}^{n}c_j2^{-j}=32-2(n+4)^22^{-n}$, whose last term tends to zero.
The remaining assertions follow from the integer values of $U_n$ and from
$c_{n+1}-c_n=4n+10\ge10$.
\end{proof}

\noindent The series is the one starting at $j=1$, whose initial tail state is
$U_0=32$; under the opening convention $\sum_{n\ge0}c_n2^{-(n+1)}$ the same
word has value $18$, and the two ranges must be kept apart.  The formal
construction checks the
\href{https://github.com/wcook04/plectis-erdos/blob/3d6d938d696fed0fb71dd55115a18a73738ff223/lean/ErdosProblems/Erdos251/PrimeGapDyadicTail.lean\#L1596}{recurrence},
\href{https://github.com/wcook04/plectis-erdos/blob/3d6d938d696fed0fb71dd55115a18a73738ff223/lean/ErdosProblems/Erdos251/PrimeGapDyadicTail.lean\#L1632}{strict
growth},
\href{https://github.com/wcook04/plectis-erdos/blob/3d6d938d696fed0fb71dd55115a18a73738ff223/lean/ErdosProblems/Erdos251/PrimeGapDyadicTail.lean\#L1654}{integrality
of every shift},
\href{https://github.com/wcook04/plectis-erdos/blob/3d6d938d696fed0fb71dd55115a18a73738ff223/lean/ErdosProblems/Erdos251/PrimeGapDyadicTail.lean\#L1644}{exclusion
of adjacent differences of size two}, and the
\href{https://github.com/wcook04/plectis-erdos/blob/3d6d938d696fed0fb71dd55115a18a73738ff223/lean/ErdosProblems/Erdos251/PolynomialGapSeriesValue.lean\#L92}{value
$32$ for the shifted series}, whose summand is $c_{n+1}/2^{\,n+1}$; the
rational value is recorded at
\href{https://github.com/wcook04/plectis-erdos/blob/3d6d938d696fed0fb71dd55115a18a73738ff223/lean/ErdosProblems/Erdos251/PolynomialGapSeriesValue.lean\#L109}{line~109}.
Positivity, parity, strict growth, unboundedness and nonperiodicity are
therefore jointly compatible with rationality, and the adjacent trigger of
Theorem~\ref{long251:res:signedwindow} fails at every index.  The telescoping mechanism
is the one used in the note of ChatGPT~5.4 Pro (orchestrated by Vjeko
Kova\v{c}) against the variable-denominator expectation of
Erd\H{o}s~\cite[Theorem~1 and proof, pp.~1--2]{kovac2026}; the word is not a
result about Problem~\#251.

Rationality also fails to force the coefficients to repeat.  Let $K:\N\to\Q$ be
arbitrary and put $\kappa_n=2K_n-K_{n+1}$: the
\href{https://github.com/wcook04/plectis-erdos/blob/3d6d938d696fed0fb71dd55115a18a73738ff223/lean/ErdosProblems/Erdos251/PrimeGapDyadicTail.lean\#L1129}{carry coefficient}.

\begin{proposition}[exact telescoping]\label{long251:res:telescope}\compatlabel{res:telescope}\compatlabel{sec:carry}
For every $n\ge0$,
$\sum_{i=0}^{n-1}\kappa_i2^{-(i+1)}=K_0-K_n2^{-n}$.
\end{proposition}

\begin{proof}
A routine induction; the added term is
$(2K_n-K_{n+1})2^{-(n+1)}=K_n2^{-n}-K_{n+1}2^{-(n+1)}$.
\end{proof}

\noindent The carry form $\kappa_n=2K_n-K_{n+1}$ is the case $a_n=2$ of the
rationality criterion of Erd\H{o}s and Straus: for integers $b_n$ and
positive integers $a_n$ with $a_n>1$ for all large $n$ and
$|b_n|/(a_{n-1}a_n)\to0$, the series $\sum_nb_n/(a_1\cdots a_n)$ is rational
exactly when some positive integer $B$ and integers $c_n$ satisfy
$Bb_n=c_na_n-c_{n+1}$ and $|c_{n+1}|<a_n/2$ for all large
$n$~\cite[Theorem~2.1, pp.~85--86]{erdosstraus1974}.  For the prime gaps the
numerators $b_n=g_n$ are unbounded, so that hypothesis fails, and the
proposition above imposes no condition on $K$.

\noindent Formalised as the
\href{https://github.com/wcook04/plectis-erdos/blob/3d6d938d696fed0fb71dd55115a18a73738ff223/lean/ErdosProblems/Erdos251/PrimeGapDyadicTail.lean\#L1137}{carry
telescoping identity}.  Taking $K_0=\tfrac52$ and $K_n=2n+2$ for $n\ge1$ gives
$\kappa_0=1$ and $\kappa_n=2n$, so the coefficients $1,2,4,6,8,\ldots$ are
positive, even after the first term, unbounded and not eventually periodic,
while their dyadic sum is $\tfrac52$.  Rationality alone therefore cannot imply
eventual periodicity even for a positive, parity-correct integer coefficient
sequence, and no argument may play periodicity of a rational value against
Proposition~\ref{long251:res:gap-nonperiodic}.

The affine and fixed-lattice escape formulations do not create new
prime-distribution information.  Their full definitions and exact
hypotheses are preserved in Appendix~\ref{long251:app:affine-detail}.

\section{What is measured}\label{long251:sec:measurements}

Two finite computations report on the two producers.  Each covers one bounded
range and supplies no cofinal statement.

Over the $6\,841\,648$ primes below $1.2\times10^{8}$ and offsets
$h=1,\ldots,16$, the reduced two-window event of
Theorem~\ref{long251:res:signedwindow} occurs at every offset, with density between
$0.00418$ and $0.008248$ and the two signs near balanced; at $h=1$ there are
$56\,427$ occurrences among $6\,841\,564$ candidate indices, the last at the
prime $119\,995\,753$.  Rescaled by $\log p$ the band densities run from
$0.17671$ down to $0.13148$ with decelerating drift over the measured bands.
This finite observation does not establish an asymptotic counting law or
cofinality.  The tails in this scan are double-precision
floating-point, so an individual hit is a numerical observation and carries no
certificate; the median distance from the shift to the nearer window boundary
is $0.127$ and the worst case over $34\,000$ hits is $6.1\times10^{-7}$, eight
orders above the working resolution, and $3\,200$ sampled hits across
$h=1,\ldots,8$ were rechecked against the literal three-part statement with no
violations.  Proposition~\ref{long251:res:finite-smallpair} is the one pair certified
by exact integer arithmetic.  The receipt is
\href{https://github.com/wcook04/plectis-erdos/blob/dd6cc708f650f061de06594eba6ba43878f4de31/research/experiments/erdos251/receipts/adjacent-mismatch.json}{\texttt{adjacent-mismatch.json}}.

Over the $1\,270\,607$ primes below $2\times10^{7}$, and for every modulus
$t\le20$ and every residue class modulo $t$, the free-pair event of
Theorem~\ref{long251:res:freepair} has witnesses running to the last usable index: the
leanest class at $t=20$ carries $90\,150$ witnesses and the latest witness sits
at index $1\,270\,520$ of $1\,270\,540$.  The recorded floating-point witnesses
pass the stronger two-window test over this finite range; they are numerical
observations, not exact certificates or a proof of either cofinal producer.  What the
scan does not supply is uniformity in $t$ and in the cutoff.  The receipt is
\href{https://github.com/wcook04/plectis-erdos/blob/dd6cc708f650f061de06594eba6ba43878f4de31/research/experiments/erdos251/receipts/free-pair.json}{\texttt{free-pair.json}}.
The \href{https://github.com/wcook04/plectis-erdos/blob/dd6cc708f650f061de06594eba6ba43878f4de31/research/experiments/erdos251/README.md}{public computation guide}
provides the programs, dependencies and exact replay commands for all three
recorded runs, including the continued-fraction calculation above.

\section{The remaining obligation}\label{long251:sec:open}

Problem~\#251 is open.  The exact unresolved condition, equivalent to
irrationality by Theorem~\ref{long251:res:escape-irrational} and the escape
equivalence for the actual tail, is the following.

\begin{problem}[universal prime-gap shift escape]\label{long251:prob:escape}\compatlabel{prob:escape}
For every $h\ge1$ and every $N_0$, prove that some $N\ge N_0$ satisfies
\begin{equation}
 \sum_{j\ge1}\frac{g_{N+h+j}-g_{N+j}}{2^{\,j}}\notin\Z .
\label{long251:eq:shift-escape}\compatlabel{eq:shift-escape}
\end{equation}
\end{problem}

\begin{problem}[cofinal adjacent small mismatch]\label{long251:prob:smallpair}
For every $h\ge1$ and every $N_0$, prove that some $N\ge N_0$ satisfies
\begin{equation}
 \bigl|\sigma_h(N)\bigr|<1,\qquad
 \bigl|\sigma_h(N+1)\bigr|<1,\qquad
 g_{N+h+1}\ne g_{N+1} .
\label{long251:eq:smallpair}
\end{equation}
\end{problem}

Corollary~\ref{long251:res:smallpair-real} makes Problem~\ref{long251:prob:smallpair}
sufficient for Problem~\ref{long251:prob:escape}, and Theorem~\ref{long251:res:freepair}
supplies an equivalent target in which the offset is free.  The three
formulations are equivalent or ordered by implication;
none of them is an analytic saving on its own.

What the proved obstructions leave.  By Theorem~\ref{long251:res:boundedperturbation}
no hypothesis stable under bounded residue-preserving perturbation can work;
by Corollary~\ref{long251:res:nonconc-primes} adding fixed-block polynomial
nonconcentration to that list does not repair it; the construction in
Section~\ref{long251:res:all-residue-log-countermodel} shows that recurring gap
values differing by $2$ in every residue class do not work at the level of
integer recurrences; and by
Theorem~\ref{long251:res:sparse} the event in \eqref{long251:eq:smallpair} has density zero, so
it must be produced cofinally on a sparse set.  Isolated small gaps, isolated
large gaps and average gap estimates are insufficient.  Any prescribed finite
prefix can be preserved while the complete sum is rationalised, so that prefix
alone cannot force irrationality or a cofinal small-tail producer.  A finite
block together with a proved tail bound can still certify one window: both
inequalities in \eqref{long251:eq:smallpair} contain the complete infinite
continuation, and Proposition~\ref{long251:res:explicit-remainder} supplies the
corresponding finite reduction.  Nor does
parity help, since after the first gap every $g_n$ is even.

The finite reduction is available.  Proposition~\ref{long251:res:explicit-remainder}
turns each window condition into an integer comparison at truncation length
$L$, and for fixed $h$ and $\varepsilon>0$ the choice
$L=\lceil(4+\varepsilon)\log_2(N+2)\rceil$ makes $E_{h,N,L}=O_h(N^{-\varepsilon})$.
A prescribed positive margin $\eta_h$ therefore converts \eqref{long251:eq:smallpair}
into the stronger task of producing the corresponding fixed-margin
finite-block events.  Strict open-window membership alone need not supply
such a fixed margin.  The finite event concerns the joint distribution
modulo powers of two of the block $(g_{N+h+1}-g_{N+1},\ldots,g_{N+h+L}-g_{N+L})$ over a window of
logarithmic length.  Deciding one instance costs $O(\log N)$ gaps together with
the elementary prime bound; producing certified witnesses beyond every cutoff, for each fixed $h$,
is the unproved input.  Uniform constants simultaneously in all $h$
are not required by the irrationality criterion.

The following published results on prime gaps address a different shape of
question.  Zhang's bounded-gap theorem~\cite[Theorem~1, p.~1122]{zhang2014}
produces infinitely many bounded consecutive-prime gaps.  Maynard bounds
$\liminf_n(p_{n+m}-p_n)$ for every fixed $m$, giving bounded clusters of every
fixed size~\cite[Theorem~1.1, p.~384]{maynard2015}, with the explicit
unconditional bound
$\liminf_n(p_{n+1}-p_n)\le600$~\cite[Theorem~1.3, p.~385]{maynard2015}.
Polymath subsequently improved this to $246$
\cite[Theorem~1.4(i) of arXiv v4]{polymath2014}; neither numerical bound
controls the signed weighted tails required here.  The large-gap
theorem of Ford, Green, Konyagin, Maynard and Tao gives an effective lower
bound for the largest single consecutive-prime gap below
$X$~\cite[Theorem~1]{fgkmt2018}.  A theorem giving bounded clusters at
each fixed cluster size does not by itself supply a weighted condition on
windows whose length grows with the basepoint, and none of these results is
used as a proof input here.

Conditionally the picture is different.  Tao's comment names uniform
quantitative prime-tuples control of about $\log\log n$ consecutive gaps as the
plausible route~\cite[comment of 7 October 2025]{erdosproblems251thread}, and
Land's draft carries that out under Kuperberg's uniform prime-tuples
conjecture~\cite{land2026}.  The obstructions of
Section~\ref{long251:sec:obstructions} say which unconditional substitutes cannot
replace that hypothesis.

\paragraph{Verification boundary.}
Ordinary proofs, exact executable certificates and checked proof
statements are distinct evidence types.  The supplied declaration index
already contains formal source for some results previously described as
having no Lean declaration, but source availability need not mean inclusion
in the checked build.  Use the per-declaration status ledger rather than
that obsolete blanket claim.  In particular, the larger
continued-fraction computation is not promoted to a kernel-decided result
merely by the existence of its verification code.  The smaller
kernel-decided denominator floor and the larger exact-arithmetic exclusion
have different evidence types and different numerical strengths.

\pdfbookmark[1]{Statements and declarations}{statements-declarations}
\subsection*{Statements and declarations}

\paragraph{Artefact and data availability.}
Declaration links identify immutable revisions, not a live branch.  The
public source revision is \texttt{3d6d938d696f}; release-only declarations
are linked to \texttt{52f29ad173b0}.  The supplied declaration index
distinguishes \texttt{ci\_checked}, \texttt{outside\_checked\_build}, and
\texttt{release\_only}; presence of a file or link alone is not evidence
that its proof was checked.  The supplied public CI receipt is
\href{https://github.com/wcook04/plectis-erdos/actions/runs/35073961520}{run 35073961520}.
The companion verification ledger records the precise declaration status.
This revision replayed \texttt{lake build ErdosProblems Erdos249257} on
Lean~4.29.1; no new \#251 endpoints were admitted beyond that replay.

Lean~4~\cite{lean4} and mathlib~\cite{mathlib} provide the proof-checking
environment.  A checked proof statement, its hypotheses and its transitive
axiom report must be distinguished from a challenge containing
\texttt{sorry}, an executable certificate, or an ordinary proof.  Formal
checking does not establish novelty, attribution, or the faithfulness of
an unexamined mathematical interpretation.

The finite-computation receipts and generating programs remain separately
identified in the
\href{https://github.com/wcook04/plectis-erdos/blob/dd6cc708f650f061de06594eba6ba43878f4de31/research/experiments/erdos251/README.md}{pinned computation guide}.
The revision packet additionally supplies an independent integer replay of
the adjacent-window certificate and the denominator exclusion
$q\ge2^{39997}$.  That replay is not a Lean check and does not authenticate
the original continued-fraction run.

\paragraph{Funding and competing interests.} This work received no external
funding.  The author declares no competing interests.

\paragraph{Acknowledgements.}
The problem numbering follows the Erd\H{o}s Problems catalogue maintained
by Thomas Bloom~\cite{erdosproblems}.  The logarithmic-scale construction
and nonconcentration transfer are retained from the earlier review
materials accompanying this project.  That provenance does not constitute
an independent specialist endorsement of the manuscripts.  No such
endorsement is claimed.  Any named acknowledgement for future correspondence
requires the contributor's permission.

\appendix
\section{An elementary polynomial bound for the primes}\label{long251:app:prime-bound}

We prove $p_n\le1250(n+1)^4$ for every $n\ge0$ by the central binomial
coefficient method of Erd\H{o}s's proof of Chebyshev's
theorem~\cite[\S1, pp.~194--196]{erdos1932}.  Let $\pi(x)$ count the primes
at most $x$.  For an integer $m\ge4$,
\begin{equation}\label{long251:eq:binomial-bound}
 4^m<m\binom{2m}{m}\le m(2m)^{\pi(2m)} .
\end{equation}
For the first inequality, $m\binom{2m}{m}/4^{m}$ equals $70/64$ at $m=4$ and
its ratio at successive indices is $(2m+1)/(2m)>1$.  For the second, the
exponent of a prime $\ell$ in $\binom{2m}{m}$ is
$\sum_{k\ge1}\bigl(\lfloor2m/\ell^{k}\rfloor-2\lfloor m/\ell^{k}\rfloor\bigr)$,
each summand is $0$ or $1$, and every summand with $\ell^{k}>2m$ vanishes;
hence each prime power dividing $\binom{2m}{m}$ is at most $2m$, and there are
$\pi(2m)$ of them.

Now fix $n\ge0$, put $x=n+5$ and $m=x^4\ge625$, and suppose $\pi(2m)\le n$.
Since $x\le2^{x}$ and $n+4(n+1)x=4x^2-15x-5\le2x^4$,
\[
 m(2m)^{n}=2^{n}x^{4(n+1)}
 \le2^{\,n+4(n+1)x}\le2^{\,2x^4}=4^{m},
\]
contradicting \eqref{long251:eq:binomial-bound}.  Hence $\pi(2m)>n$, so at least $n+1$
primes lie below $2m$ and therefore
\[
 p_n\le2(n+5)^4\le1250(n+1)^4 ,
\]
the last step because $(n+5)\le5(n+1)$ for $n\ge0$.  The same bound is checked
by the kernel at
\href{https://github.com/wcook04/plectis-erdos/blob/3d6d938d696fed0fb71dd55115a18a73738ff223/lean/ErdosProblems/Erdos251/PrimeGapDyadicTail.lean\#L360}{line~360}.

\section{Integer certificates}\label{long251:app:certificates}

The following calculation uses trial division and integer arithmetic only.  It
reproduces the two rows of Proposition~\ref{long251:res:finite-smallpair} and every
inequality used in Theorem~\ref{long251:res:denominatorfloor}.  The four literals are
the certificate the Lean kernel decides in
\texttt{KernelDenominatorFloor.lean}; their use here requires no
floating-point approximation.  Adjacent quoted strings are concatenated by
Python.
{\footnotesize
\begin{verbatim}
from math import isqrt

primes = [n for n in range(2, 10000)
          if all(n % d for d in range(2, isqrt(n) + 1))]
gaps = [q - p for p, q in zip(primes, primes[1:])]
P = lambda x: x**4 + 8*x**3 + 36*x*x + 104*x + 150
Q = 2**40
expected = [(-662838684750, 11764181250),
            (873345886050, 12805761250)]
for N, target in zip((2, 3), expected):
    D = sum((gaps[N+1+j] - gaps[N+j])*2**(40-j)
            for j in range(1, 41))
    B = 1250*(P(N+43) + P(N+42))
    distance = min(D % Q, Q - D % Q)
    if (D, B) != target or not (abs(D)+B < Q and B < distance):
        raise ArithmeticError("tail certificate failed")
if gaps[4] == gaps[3]:
    raise ArithmeticError("gap mismatch failed")

u = int(
    "8065641857152652932176019632186898003271162829171466334827"
    "3083607794415278717445033509407855988903369988525550746159"
    "7355889792250084202344821020139160956663658789718168152662"
    "0217"
)

v = int(
    "2194945124413663232143970924541263312422069524635615360518"
    "4247077351958221816830720189289904831662955084392690248683"
    "1216291723988537733235173040607254496838530213867781442335"
    "1745"
)

up = int(
    "6539437101498162626882411892470905222108260008568555445975"
    "3026163315521784668689909712759876562484620059098138417469"
    "5232839888185316374272333277389611483117334000493867584923"
    "912"
)

vp = int(
    "1779611075789866551198427241621709630120136323281598176637"
    "8490605249229735501968764904036152970729491656650873938580"
    "6203025466313389810291243786005499164537499383612081805160"
    "767"
)

c = len(primes)
A = sum(p*2**(c-1-i) for i, p in enumerate(primes))
checks = (c == 1229,
          u > 0 and v > 0 and up > 0 and vp > 0,
          up*v - u*vp == 1,
          u*2**c < A*v,
          (2*A + 5000*(c+1)**4)*vp < up*2**(c+1),
          v+vp >= 2**589,
          2**589 > 10**177)
if not all(checks):
    raise ArithmeticError("denominator certificate failed")
print("Both tail rows and the denominator floor verified.")
\end{verbatim}
}

\section{Guide to the formal sources}\label{long251:app:index}\compatlabel{app:index}

Each link identifies its own immutable source revision.  Named declaration
links have been refreshed from the supplied checked index; historical
computation links retain their separate receipt pins.  The summation-by-parts and prime-bound declarations
are prime-specific.  Section~\ref{long251:sec:tail} is stated for arbitrary integer
digits and arbitrary rational or real orbits, while
Section~\ref{long251:sec:local-certificate} records the actual-gap specialisation.  The
concrete prime-gap tail, its unconditional convergence, its recurrence and the
real-to-rational scaled-tail bridge are defined and checked in
\texttt{RealPrimeGapTail.lean} and \texttt{PrimeGapDyadicTail.lean}; the
free-pair lattice and its actual-series equivalence in
\texttt{FreePairReduction.lean}; the kernel-decided denominator certificate in
\texttt{KernelDenominatorFloor.lean}; the rational countermodels in
\texttt{BoundedPerturbationCountermodel.lean} and
\texttt{PolynomialGapSeriesValue.lean}; and the two circularity theorems in
\texttt{AffineCylinderCollapse.lean}, with the signed normal form in
\texttt{AffineShiftEscape.lean} and the lcm diagonal in
\texttt{OrderLatticeDiagonal.lean}.

% The exhaustive source manifest is retained in the authored TeX for automated
% validation, and kept outside the printed note.
\iffalse

\begin{tabular}{@{}ll@{}}
informal statement & linked Lean source\\
zero-based primes & \href{https://github.com/wcook04/plectis-erdos/blob/3d6d938d696fed0fb71dd55115a18a73738ff223/lean/ErdosProblems/Erdos251/PrimeGapDyadicTail.lean\#L40}{\leanlabel{prime0}}\\
zero-based gaps & \href{https://github.com/wcook04/plectis-erdos/blob/3d6d938d696fed0fb71dd55115a18a73738ff223/lean/ErdosProblems/Erdos251/PrimeGapDyadicTail.lean\#L44}{\leanlabel{primeGap0}}\\
first gap value & \href{https://github.com/wcook04/plectis-erdos/blob/3d6d938d696fed0fb71dd55115a18a73738ff223/lean/ErdosProblems/Erdos251/PrimeGapDyadicTail.lean\#L47}{\leanlabel{primeGap0_zero}}\\
second gap value & \href{https://github.com/wcook04/plectis-erdos/blob/3d6d938d696fed0fb71dd55115a18a73738ff223/lean/ErdosProblems/Erdos251/PrimeGapDyadicTail.lean\#L51}{\leanlabel{primeGap0_one}}\\
unbounded actual prime gaps & \href{https://github.com/wcook04/plectis-erdos/blob/3d6d938d696fed0fb71dd55115a18a73738ff223/lean/ErdosProblems/Erdos251/PrimeGapDyadicTail.lean\#L57}{\leanlabel{exists_primeGap0_gt}}\\
dyadic partial sum & \href{https://github.com/wcook04/plectis-erdos/blob/3d6d938d696fed0fb71dd55115a18a73738ff223/lean/ErdosProblems/Erdos251/PrimeGapDyadicTail.lean\#L101}{\leanlabel{dyadicPartialSumQ}}\\
displayed normalisation & \href{https://github.com/wcook04/plectis-erdos/blob/3d6d938d696fed0fb71dd55115a18a73738ff223/lean/ErdosProblems/Erdos251/PrimeGapDyadicTail.lean\#L106}{\leanlabel{prime0DisplayedPartialSumQ}}\\
factor-of-two identity & \href{https://github.com/wcook04/plectis-erdos/blob/3d6d938d696fed0fb71dd55115a18a73738ff223/lean/ErdosProblems/Erdos251/PrimeGapDyadicTail.lean\#L111}{\leanlabel{prime0DisplayedPartialSumQ_eq_two_mul}}\\
dyadic difference sum & \href{https://github.com/wcook04/plectis-erdos/blob/3d6d938d696fed0fb71dd55115a18a73738ff223/lean/ErdosProblems/Erdos251/PrimeGapDyadicTail.lean\#L121}{\leanlabel{dyadicDifferencePartialSumQ}}\\
partial-sum step & \href{https://github.com/wcook04/plectis-erdos/blob/3d6d938d696fed0fb71dd55115a18a73738ff223/lean/ErdosProblems/Erdos251/PrimeGapDyadicTail.lean\#L124}{\leanlabel{dyadicPartialSumQ_succ}}\\
difference-sum step & \href{https://github.com/wcook04/plectis-erdos/blob/3d6d938d696fed0fb71dd55115a18a73738ff223/lean/ErdosProblems/Erdos251/PrimeGapDyadicTail.lean\#L129}{\leanlabel{dyadicDifferencePartialSumQ_succ}}\\
summation by parts & \href{https://github.com/wcook04/plectis-erdos/blob/3d6d938d696fed0fb71dd55115a18a73738ff223/lean/ErdosProblems/Erdos251/PrimeGapDyadicTail.lean\#L138}{\leanlabel{dyadicPartialSumQ_eq_start_add_differences}}\\
monotonicity step & \href{https://github.com/wcook04/plectis-erdos/blob/3d6d938d696fed0fb71dd55115a18a73738ff223/lean/ErdosProblems/Erdos251/PrimeGapDyadicTail.lean\#L152}{\leanlabel{prime0_mono_step}}\\
gap cast identity & \href{https://github.com/wcook04/plectis-erdos/blob/3d6d938d696fed0fb71dd55115a18a73738ff223/lean/ErdosProblems/Erdos251/PrimeGapDyadicTail.lean\#L156}{\leanlabel{primeGap0_cast}}\\
gap partial sum & \href{https://github.com/wcook04/plectis-erdos/blob/3d6d938d696fed0fb71dd55115a18a73738ff223/lean/ErdosProblems/Erdos251/PrimeGapDyadicTail.lean\#L161}{\leanlabel{primeGapPartialSumQ}}\\
prime-gap summation by parts & \href{https://github.com/wcook04/plectis-erdos/blob/3d6d938d696fed0fb71dd55115a18a73738ff223/lean/ErdosProblems/Erdos251/PrimeGapDyadicTail.lean\#L172}{\leanlabel{prime0_dyadic_summation_by_parts}}\\
prime dyadic term & \href{https://github.com/wcook04/plectis-erdos/blob/3d6d938d696fed0fb71dd55115a18a73738ff223/lean/ErdosProblems/Erdos251/PrimeGapDyadicTail.lean\#L183}{\leanlabel{primeDyadicTerm}}\\
displayed prime dyadic term & \href{https://github.com/wcook04/plectis-erdos/blob/3d6d938d696fed0fb71dd55115a18a73738ff223/lean/ErdosProblems/Erdos251/PrimeGapDyadicTail.lean\#L188}{\leanlabel{primeDisplayedDyadicTerm}}\\
prime-gap dyadic term & \href{https://github.com/wcook04/plectis-erdos/blob/3d6d938d696fed0fb71dd55115a18a73738ff223/lean/ErdosProblems/Erdos251/PrimeGapDyadicTail.lean\#L192}{\leanlabel{primeGapDyadicTerm}}\\
infinite factor-of-two identity & \href{https://github.com/wcook04/plectis-erdos/blob/3d6d938d696fed0fb71dd55115a18a73738ff223/lean/ErdosProblems/Erdos251/PrimeGapDyadicTail.lean\#L196}{\leanlabel{primeDisplayedDyadicTerm_eq_two_mul}}\\
gap discrete derivative & \href{https://github.com/wcook04/plectis-erdos/blob/3d6d938d696fed0fb71dd55115a18a73738ff223/lean/ErdosProblems/Erdos251/PrimeGapDyadicTail.lean\#L202}{\leanlabel{primeGapDyadicTerm_eq}}\\
polynomial prime bound & \href{https://github.com/wcook04/plectis-erdos/blob/3d6d938d696fed0fb71dd55115a18a73738ff223/lean/ErdosProblems/Erdos251/PrimeGapDyadicTail.lean\#L360}{\leanlabel{prime0_le_polynomial}}\\
prime-series summability & \href{https://github.com/wcook04/plectis-erdos/blob/3d6d938d696fed0fb71dd55115a18a73738ff223/lean/ErdosProblems/Erdos251/PrimeGapDyadicTail.lean\#L379}{\leanlabel{summable_primeDyadicTerm}}\\
gap summability transfer & \href{https://github.com/wcook04/plectis-erdos/blob/3d6d938d696fed0fb71dd55115a18a73738ff223/lean/ErdosProblems/Erdos251/PrimeGapDyadicTail.lean\#L385}{\leanlabel{summable_primeGapDyadicTerm_of_summable_primeDyadicTerm}}\\
infinite gap identity & \href{https://github.com/wcook04/plectis-erdos/blob/3d6d938d696fed0fb71dd55115a18a73738ff223/lean/ErdosProblems/Erdos251/PrimeGapDyadicTail.lean\#L404}{\leanlabel{tsum_primeDyadicTerm_eq_two_add_primeGap}}\\
unconditional gap identity & \href{https://github.com/wcook04/plectis-erdos/blob/3d6d938d696fed0fb71dd55115a18a73738ff223/lean/ErdosProblems/Erdos251/PrimeGapDyadicTail.lean\#L427}{\leanlabel{tsum_primeDyadicTerm_eq_two_add_primeGap_unconditional}}\\
normalised irrationality equivalence & \href{https://github.com/wcook04/plectis-erdos/blob/3d6d938d696fed0fb71dd55115a18a73738ff223/lean/ErdosProblems/Erdos251/PrimeGapDyadicTail.lean\#L435}{\leanlabel{irrational_tsum_primeDyadicTerm_iff_primeGap}}\\
displayed infinite identity & \href{https://github.com/wcook04/plectis-erdos/blob/3d6d938d696fed0fb71dd55115a18a73738ff223/lean/ErdosProblems/Erdos251/PrimeGapDyadicTail.lean\#L444}{\leanlabel{tsum_primeDisplayedDyadicTerm_eq_four_add_two_primeGap}}\\
displayed irrationality equivalence & \href{https://github.com/wcook04/plectis-erdos/blob/3d6d938d696fed0fb71dd55115a18a73738ff223/lean/ErdosProblems/Erdos251/PrimeGapDyadicTail.lean\#L459}{\leanlabel{irrational_tsum_primeDisplayedDyadicTerm_iff_primeGap}}\\
tail recurrence & \href{https://github.com/wcook04/plectis-erdos/blob/3d6d938d696fed0fb71dd55115a18a73738ff223/lean/ErdosProblems/Erdos251/PrimeGapDyadicTail.lean\#L482}{\leanlabel{DyadicTailRecurrence}}\\
rational actual-tail state & \href{https://github.com/wcook04/plectis-erdos/blob/3d6d938d696fed0fb71dd55115a18a73738ff223/lean/ErdosProblems/Erdos251/PrimeGapDyadicTail.lean\#L489}{\leanlabel{rationalPrimeGapTailState}}\\
hypothetical-rational tail representation & \href{https://github.com/wcook04/plectis-erdos/blob/3d6d938d696fed0fb71dd55115a18a73738ff223/lean/ErdosProblems/Erdos251/PrimeGapDyadicTail.lean\#L510}{\leanlabel{exists_rationalPrimeGapTailState_representation_of_not_irrational}}\\
rational actual-tail recurrence & \href{https://github.com/wcook04/plectis-erdos/blob/3d6d938d696fed0fb71dd55115a18a73738ff223/lean/ErdosProblems/Erdos251/PrimeGapDyadicTail.lean\#L527}{\leanlabel{rationalPrimeGapTailState_recurrence}}\\
shift & \href{https://github.com/wcook04/plectis-erdos/blob/3d6d938d696fed0fb71dd55115a18a73738ff223/lean/ErdosProblems/Erdos251/PrimeGapDyadicTail.lean\#L544}{\leanlabel{tailShift}}\\
shift step identity & \href{https://github.com/wcook04/plectis-erdos/blob/3d6d938d696fed0fb71dd55115a18a73738ff223/lean/ErdosProblems/Erdos251/PrimeGapDyadicTail.lean\#L563}{\leanlabel{tailShift_succ}}\\
integrality & \href{https://github.com/wcook04/plectis-erdos/blob/3d6d938d696fed0fb71dd55115a18a73738ff223/lean/ErdosProblems/Erdos251/PrimeGapDyadicTail.lean\#L575}{\leanlabel{RatIntegral}}\\
tail block & \href{https://github.com/wcook04/plectis-erdos/blob/3d6d938d696fed0fb71dd55115a18a73738ff223/lean/ErdosProblems/Erdos251/PrimeGapDyadicTail.lean\#L647}{\leanlabel{dyadicTailBlock}}\\
iterated block identity & \href{https://github.com/wcook04/plectis-erdos/blob/3d6d938d696fed0fb71dd55115a18a73738ff223/lean/ErdosProblems/Erdos251/PrimeGapDyadicTail.lean\#L653}{\leanlabel{tail_iterate_eq_pow_mul_sub_block}}\\
scaled shift identity & \href{https://github.com/wcook04/plectis-erdos/blob/3d6d938d696fed0fb71dd55115a18a73738ff223/lean/ErdosProblems/Erdos251/PrimeGapDyadicTail.lean\#L667}{\leanlabel{tailShift_eq_scaled_sub_block}}\\
integer-shift invariance & \href{https://github.com/wcook04/plectis-erdos/blob/3d6d938d696fed0fb71dd55115a18a73738ff223/lean/ErdosProblems/Erdos251/PrimeGapDyadicTail.lean\#L677}{\leanlabel{ratIntegral_sub_int_iff}}\\
Euler multiplier & \href{https://github.com/wcook04/plectis-erdos/blob/3d6d938d696fed0fb71dd55115a18a73738ff223/lean/ErdosProblems/Erdos251/PrimeGapDyadicTail.lean\#L695}{\leanlabel{ratIntegral_totientMultiplier_of_odd_den}}\\
power-of-two/odd denominator scaling & \href{https://github.com/wcook04/plectis-erdos/blob/3d6d938d696fed0fb71dd55115a18a73738ff223/lean/ErdosProblems/Erdos251/PrimeGapDyadicTail.lean\#L723}{\leanlabel{ratIntegral_scaled_of_den_eq_pow_two_mul_odd}}\\
multiplier at the two-adic exponent & \href{https://github.com/wcook04/plectis-erdos/blob/3d6d938d696fed0fb71dd55115a18a73738ff223/lean/ErdosProblems/Erdos251/PrimeGapDyadicTail.lean\#L756}{\leanlabel{ratIntegral_multiplier_tail_at_twoExponent}}\\
integral-shift criterion & \href{https://github.com/wcook04/plectis-erdos/blob/3d6d938d696fed0fb71dd55115a18a73738ff223/lean/ErdosProblems/Erdos251/PrimeGapDyadicTail.lean\#L802}{\leanlabel{tailShift_integral_iff_scaledTail}}\\
fixed-denominator eventual shift & \href{https://github.com/wcook04/plectis-erdos/blob/3d6d938d696fed0fb71dd55115a18a73738ff223/lean/ErdosProblems/Erdos251/PrimeGapDyadicTail.lean\#L830}{\leanlabel{rationalPrimeGapTailShift_eventuallyIntegral_of_den_eq}}\\
eventual integral shift & \href{https://github.com/wcook04/plectis-erdos/blob/3d6d938d696fed0fb71dd55115a18a73738ff223/lean/ErdosProblems/Erdos251/PrimeGapDyadicTail.lean\#L853}{\leanlabel{rationalPrimeGapTailShift_eventuallyIntegral}}\\
totient shift & \href{https://github.com/wcook04/plectis-erdos/blob/3d6d938d696fed0fb71dd55115a18a73738ff223/lean/ErdosProblems/Erdos251/PrimeGapDyadicTail.lean\#L877}{\leanlabel{tailShift_integral_totient_of_odd_den}}\\
one-step propagation & \href{https://github.com/wcook04/plectis-erdos/blob/3d6d938d696fed0fb71dd55115a18a73738ff223/lean/ErdosProblems/Erdos251/PrimeGapDyadicTail.lean\#L888}{\leanlabel{tailShift_integral_succ}}\\
propagation to every later index & \href{https://github.com/wcook04/plectis-erdos/blob/3d6d938d696fed0fb71dd55115a18a73738ff223/lean/ErdosProblems/Erdos251/PrimeGapDyadicTail.lean\#L900}{\leanlabel{tailShift_integral_add}}\\
adjacent small-shift obstruction & \href{https://github.com/wcook04/plectis-erdos/blob/3d6d938d696fed0fb71dd55115a18a73738ff223/lean/ErdosProblems/Erdos251/PrimeGapDyadicTail.lean\#L979}{\leanlabel{tailShift_not_both_integral_of_small_pair_of_digit_ne}}\\
cofinal small-mismatch theorem & \href{https://github.com/wcook04/plectis-erdos/blob/3d6d938d696fed0fb71dd55115a18a73738ff223/lean/ErdosProblems/Erdos251/PrimeGapDyadicTail.lean\#L1006}{\leanlabel{not_eventuallyIntegralTailShift_of_cofinal_small_mismatch}}\\
prime gaps not eventually periodic & \href{https://github.com/wcook04/plectis-erdos/blob/3d6d938d696fed0fb71dd55115a18a73738ff223/lean/ErdosProblems/Erdos251/PrimeGapDyadicTail.lean\#L1023}{\leanlabel{primeGap0_not_eventually_periodic}}\\
shift not eventually small & \href{https://github.com/wcook04/plectis-erdos/blob/3d6d938d696fed0fb71dd55115a18a73738ff223/lean/ErdosProblems/Erdos251/PrimeGapDyadicTail.lean\#L1097}{\leanlabel{rationalPrimeGapTail_has_positive_shift_not_eventually_small}}\\
actual-prime-gap specialisation & \href{https://github.com/wcook04/plectis-erdos/blob/3d6d938d696fed0fb71dd55115a18a73738ff223/lean/ErdosProblems/Erdos251/PrimeGapDyadicTail.lean\#L1112}{\leanlabel{primeGapTailShift_not_eventuallyIntegral_of_cofinal_small_mismatch}}\\
carry coefficient & \href{https://github.com/wcook04/plectis-erdos/blob/3d6d938d696fed0fb71dd55115a18a73738ff223/lean/ErdosProblems/Erdos251/PrimeGapDyadicTail.lean\#L1129}{\leanlabel{carryCoeff}}\\
carry partial sum & \href{https://github.com/wcook04/plectis-erdos/blob/3d6d938d696fed0fb71dd55115a18a73738ff223/lean/ErdosProblems/Erdos251/PrimeGapDyadicTail.lean\#L1133}{\leanlabel{carryPartialSum}}\\
carry telescoping identity & \href{https://github.com/wcook04/plectis-erdos/blob/3d6d938d696fed0fb71dd55115a18a73738ff223/lean/ErdosProblems/Erdos251/PrimeGapDyadicTail.lean\#L1137}{\leanlabel{carryPartialSum_eq}}\\
natural carry coefficient & \href{https://github.com/wcook04/plectis-erdos/blob/3d6d938d696fed0fb71dd55115a18a73738ff223/lean/ErdosProblems/Erdos251/PrimeGapDyadicTail.lean\#L1177}{\leanlabel{natCarryCoeff}}\\
natural-carry cast & \href{https://github.com/wcook04/plectis-erdos/blob/3d6d938d696fed0fb71dd55115a18a73738ff223/lean/ErdosProblems/Erdos251/PrimeGapDyadicTail.lean\#L1180}{\leanlabel{natCarryCoeff_cast}}\\
exact denominator recurrence & \href{https://github.com/wcook04/plectis-erdos/blob/3d6d938d696fed0fb71dd55115a18a73738ff223/lean/ErdosProblems/Erdos251/PrimeGapDyadicTail.lean\#L1217}{\leanlabel{tail_den_succ}}\\
odd-denominator persistence & \href{https://github.com/wcook04/plectis-erdos/blob/3d6d938d696fed0fb71dd55115a18a73738ff223/lean/ErdosProblems/Erdos251/PrimeGapDyadicTail.lean\#L1226}{\leanlabel{tail_den_succ_eq_of_odd}}\\
even-denominator halving & \href{https://github.com/wcook04/plectis-erdos/blob/3d6d938d696fed0fb71dd55115a18a73738ff223/lean/ErdosProblems/Erdos251/PrimeGapDyadicTail.lean\#L1236}{\leanlabel{tail_den_succ_eq_half_of_even}}\\
integral shifts by denominator & \href{https://github.com/wcook04/plectis-erdos/blob/3d6d938d696fed0fb71dd55115a18a73738ff223/lean/ErdosProblems/Erdos251/PrimeGapDyadicTail.lean\#L1279}{\leanlabel{tailShift_integral_iff_den_dvd_mersenne}}\\
power-of-two congruence & \href{https://github.com/wcook04/plectis-erdos/blob/3d6d938d696fed0fb71dd55115a18a73738ff223/lean/ErdosProblems/Erdos251/PrimeGapDyadicTail.lean\#L1291}{\leanlabel{tailShift_integral_iff_two_pow_modEq_one}}\\
one-shift rationality classifier & \href{https://github.com/wcook04/plectis-erdos/blob/3d6d938d696fed0fb71dd55115a18a73738ff223/lean/ErdosProblems/Erdos251/PrimeGapDyadicTail.lean\#L1479}{\leanlabel{not_irrational_initial_iff_exists_integral_positive_tailShift}}\\
eventual-shift rationality classifier & \href{https://github.com/wcook04/plectis-erdos/blob/3d6d938d696fed0fb71dd55115a18a73738ff223/lean/ErdosProblems/Erdos251/PrimeGapDyadicTail.lean\#L1525}{\leanlabel{not_irrational_initial_iff_exists_eventually_integral_positive_tailShift}}\\
pointwise irrationality classifier & \href{https://github.com/wcook04/plectis-erdos/blob/3d6d938d696fed0fb71dd55115a18a73738ff223/lean/ErdosProblems/Erdos251/PrimeGapDyadicTail.lean\#L1551}{\leanlabel{irrational_initial_iff_all_positive_tailShifts_nonintegral}}\\
cofinal irrationality classifier & \href{https://github.com/wcook04/plectis-erdos/blob/3d6d938d696fed0fb71dd55115a18a73738ff223/lean/ErdosProblems/Erdos251/PrimeGapDyadicTail.lean\#L1572}{\leanlabel{irrational_initial_iff_cofinalNonintegralTailShifts}}\\
polynomial countermodel recurrence & \href{https://github.com/wcook04/plectis-erdos/blob/3d6d938d696fed0fb71dd55115a18a73738ff223/lean/ErdosProblems/Erdos251/PrimeGapDyadicTail.lean\#L1596}{\leanlabel{polynomialTailOrbit_recurrence}}\\
polynomial countermodel strict growth & \href{https://github.com/wcook04/plectis-erdos/blob/3d6d938d696fed0fb71dd55115a18a73738ff223/lean/ErdosProblems/Erdos251/PrimeGapDyadicTail.lean\#L1632}{\leanlabel{polynomialGapWord_strictMono}}\\
polynomial countermodel excludes the adjacent $\pm2$ trigger & \href{https://github.com/wcook04/plectis-erdos/blob/3d6d938d696fed0fb71dd55115a18a73738ff223/lean/ErdosProblems/Erdos251/PrimeGapDyadicTail.lean\#L1644}{\leanlabel{polynomialGapWord_no_adjacent_two_digit}}\\
all polynomial-countermodel shifts are integral & \href{https://github.com/wcook04/plectis-erdos/blob/3d6d938d696fed0fb71dd55115a18a73738ff223/lean/ErdosProblems/Erdos251/PrimeGapDyadicTail.lean\#L1654}{\leanlabel{polynomialTailOrbit_shift_integral}}\\
polynomial countermodel dyadic value & \href{https://github.com/wcook04/plectis-erdos/blob/3d6d938d696fed0fb71dd55115a18a73738ff223/lean/ErdosProblems/Erdos251/PolynomialGapSeriesValue.lean\#L92}{\leanlabel{tsum_polynomialGapDyadicTerm_eq}}\\
polynomial countermodel rationality & \href{https://github.com/wcook04/plectis-erdos/blob/3d6d938d696fed0fb71dd55115a18a73738ff223/lean/ErdosProblems/Erdos251/PolynomialGapSeriesValue.lean\#L109}{\leanlabel{not_irrational_tsum_polynomialGapDyadicTerm}}\\
real tail shifted-gap identity & \href{https://github.com/wcook04/plectis-erdos/blob/3d6d938d696fed0fb71dd55115a18a73738ff223/lean/ErdosProblems/Erdos251/RealPrimeGapTail.lean\#L31}{\leanlabel{realPrimeGapTail_eq_tsum_shifted_gaps}}\\
real tail recurrence & \href{https://github.com/wcook04/plectis-erdos/blob/3d6d938d696fed0fb71dd55115a18a73738ff223/lean/ErdosProblems/Erdos251/RealPrimeGapTail.lean\#L48}{\leanlabel{realPrimeGapTail_recurrence}}\\
actual escape equivalence & \href{https://github.com/wcook04/plectis-erdos/blob/3d6d938d696fed0fb71dd55115a18a73738ff223/lean/ErdosProblems/Erdos251/RealPrimeGapTail.lean\#L73}{\leanlabel{irrational_primeSeries_iff_realPrimeGapTail_escape}}\\
real adjacent-pair consumer & \href{https://github.com/wcook04/plectis-erdos/blob/3d6d938d696fed0fb71dd55115a18a73738ff223/lean/ErdosProblems/Erdos251/RealPrimeGapTail.lean\#L99}{\leanlabel{realTailShift_not_both_integral_of_small_pair_of_digit_ne}}\\
real small-mismatch criterion & \href{https://github.com/wcook04/plectis-erdos/blob/3d6d938d696fed0fb71dd55115a18a73738ff223/lean/ErdosProblems/Erdos251/RealPrimeGapTail.lean\#L120}{\leanlabel{irrational_primeSeries_of_realPrimeGapTail_small_mismatch}}\\
free-pair lattice & \href{https://github.com/wcook04/plectis-erdos/blob/3d6d938d696fed0fb71dd55115a18a73738ff223/lean/ErdosProblems/Erdos251/FreePairReduction.lean\#L92}{\leanlabel{free_pair_integral_iff_modEq}}\\
free-pair lattice corollary & \href{https://github.com/wcook04/plectis-erdos/blob/3d6d938d696fed0fb71dd55115a18a73738ff223/lean/ErdosProblems/Erdos251/FreePairReduction.lean\#L105}{\leanlabel{exists_free_pair_lattice}}\\
actual real tail recurrence & \href{https://github.com/wcook04/plectis-erdos/blob/3d6d938d696fed0fb71dd55115a18a73738ff223/lean/ErdosProblems/Erdos251/FreePairReduction.lean\#L186}{\leanlabel{primeGapRealTail_recurrence}}\\
free-pair equivalence for the actual series & \href{https://github.com/wcook04/plectis-erdos/blob/3d6d938d696fed0fb71dd55115a18a73738ff223/lean/ErdosProblems/Erdos251/FreePairReduction.lean\#L236}{\leanlabel{irrational_primeGap_tsum_iff_cofinalFreePairNonintegral}}\\
kernel-decided certificate & \href{https://github.com/wcook04/plectis-erdos/blob/3d6d938d696fed0fb71dd55115a18a73738ff223/lean/ErdosProblems/Erdos251/KernelDenominatorFloor.lean\#L310}{\leanlabel{cert_10000}}\\
denominator floor, prime series & \href{https://github.com/wcook04/plectis-erdos/blob/3d6d938d696fed0fb71dd55115a18a73738ff223/lean/ErdosProblems/Erdos251/KernelDenominatorFloor.lean\#L316}{\leanlabel{kernel_denominator_floor}}\\
denominator floor, gap series & \href{https://github.com/wcook04/plectis-erdos/blob/3d6d938d696fed0fb71dd55115a18a73738ff223/lean/ErdosProblems/Erdos251/KernelDenominatorFloor.lean\#L324}{\leanlabel{kernel_denominator_floor_primeGap}}\\
bounded-perturbation construction & \href{https://github.com/wcook04/plectis-erdos/blob/3d6d938d696fed0fb71dd55115a18a73738ff223/lean/ErdosProblems/Erdos251/BoundedPerturbationCountermodel.lean\#L193}{\leanlabel{exists_bounded_perturbation}}\\
binary-digit encoding & \href{https://github.com/wcook04/plectis-erdos/blob/3d6d938d696fed0fb71dd55115a18a73738ff223/lean/ErdosProblems/Erdos251/BoundedPerturbationCountermodel.lean\#L134}{\leanlabel{exists_bounded_perturbation_digits}}\\
prime-anchored rational perturbation & \href{https://github.com/wcook04/plectis-erdos/blob/3d6d938d696fed0fb71dd55115a18a73738ff223/lean/ErdosProblems/Erdos251/BoundedPerturbationCountermodel.lean\#L208}{\leanlabel{exists_rational_bounded_perturbation_primeGap}}\\
perturbed prime-gap series is rational & \href{https://github.com/wcook04/plectis-erdos/blob/3d6d938d696fed0fb71dd55115a18a73738ff223/lean/ErdosProblems/Erdos251/BoundedPerturbationCountermodel.lean\#L224}{\leanlabel{not_irrational_bounded_perturbation_primeGap}}\\
perturbed gap bounds & \href{https://github.com/wcook04/plectis-erdos/blob/3d6d938d696fed0fb71dd55115a18a73738ff223/lean/ErdosProblems/Erdos251/BoundedPerturbationCountermodel.lean\#L234}{\leanlabel{perturbed_gap_bounds}}\\
signed two-window normal form & \href{https://github.com/wcook04/plectis-erdos/blob/3d6d938d696fed0fb71dd55115a18a73738ff223/lean/ErdosProblems/Erdos251/AffineShiftEscape.lean\#L113}{\leanlabel{adjacent_small_mismatch_iff_signed_two_window}}\\
pointwise affine identity & \href{https://github.com/wcook04/plectis-erdos/blob/3d6d938d696fed0fb71dd55115a18a73738ff223/lean/ErdosProblems/Erdos251/AffineCylinderCollapse.lean\#L91}{\leanlabel{ratAffinePowTwo_iff_evenIntegral}}\\
cofinal affine circularity & \href{https://github.com/wcook04/plectis-erdos/blob/3d6d938d696fed0fb71dd55115a18a73738ff223/lean/ErdosProblems/Erdos251/AffineCylinderCollapse.lean\#L162}{\leanlabel{cofinal_affinePowTwo_escape_iff_not_eventuallyIntegral}}\\
fixed-lattice circularity & \href{https://github.com/wcook04/plectis-erdos/blob/3d6d938d696fed0fb71dd55115a18a73738ff223/lean/ErdosProblems/Erdos251/AffineCylinderCollapse.lean\#L308}{\leanlabel{cofinal_blockResidue_escape_iff_not_eventuallyIntegral}}\\
lcm-diagonal criterion & \href{https://github.com/wcook04/plectis-erdos/blob/3d6d938d696fed0fb71dd55115a18a73738ff223/lean/ErdosProblems/Erdos251/OrderLatticeDiagonal.lean\#L153}{\leanlabel{irrational_initial_iff_all_lcmDiagonal_nonintegral}}\\
\end{tabular}
\fi
% ---- part extended_record ----
\clearpage
\section{Extended record: displaced results, ordinary proofs and corrections}
\label{long251:sec:erdos-251-extended-record}

This section holds mathematics that belongs to the complete record for
Problem~\#251 and sits outside the short note.  Every statement here is
labelled by its evidence class in the same vocabulary the note uses.

\subsection{The totient witness and forward propagation}
\label{long251:sec:xr-totient}

The note states the exact denominator law \eqref{long251:eq:den-law} and uses it
directly.  The classical witness that produced it is worth recording, because
it is the form in which the shift length first appears.

\begin{proposition}[a shift of totient length]\label{long251:xr:totient}\compatlabel{res:totient}
Let $T:\N\to\Q$ satisfy the dyadic tail recurrence with integer digits.  If the
reduced denominator $d$ of $T_N$ is odd, then $\sigma_{\ph(d)}(N)$ is an
integer.
\end{proposition}

\begin{proof}
Since $d$ is odd, $2$ and $d$ are coprime, so Euler's congruence gives
$d\mid2^{\ph(d)}-1$.  Writing $2^{\ph(d)}-1=dk$ and $T_N=u/d$ in lowest terms,
$(2^{\ph(d)}-1)T_N=ku$ is an integer, and the integral-shift criterion
transfers this to the shift.
\end{proof}

\begin{proposition}[propagation]\label{long251:xr:propagate}\compatlabel{res:propagate}
Let $T:\N\to\R$ satisfy $T_{n+1}=2T_n-a_{n+1}$ with integer digits,
and define $\sigma_h(N)=T_{N+h}-T_N$.  For fixed $h,N\ge0$,
if $\sigma_h(N)$ is an integer, then $\sigma_h(N+k)$ is an integer
for every $k\ge0$.
\end{proposition}

\begin{proof}
The shift step identity gives
$\sigma_h(N+1)=2\sigma_h(N)-(g_{N+h+1}-g_{N+1})$, an integer combination of an
integer and two digits.  Induct on $k$.
\end{proof}

\noindent These are the totient shift and the propagation theorems cited in the
note.  Two worked instances.  If $\operatorname{den}T_N=3$ then $\ph(3)=2$ and
$3T_N$ is an integer, so $\sigma_2(N)$ is integral while $\sigma_1(N)$ is not;
if $\operatorname{den}T_N=5$ then $\sigma_4(N)$ is integral.  In the orbit with
every digit zero and $T_0=1/12$, the denominator is $12=2^2\cdot3$, the orbit
reaches $T_2=1/3$ with odd denominator at $s=2$, and $h=\ph(3)=2$ is exactly the
shift length seen to be integral from index $2$ onwards.  The totient is a
witness rather than the classification: the exact criterion is
$\operatorname{den}(T_N)\mid2^{h}-1$, and the least admissible $h$ is the
multiplicative order of $2$ modulo the odd part.

\subsection{The finite truncation criterion}\label{long251:sec:xr-truncation}

The note gives the explicit remainder bound for the actual gaps.  The general
truncation statement behind it, valid for any dominating majorant, is the
following.

\begin{proposition}[finite truncation]\label{long251:xr:truncation}
Let $M(n)\ge g_n$ for every $n$, assume the series below converges, and put
\[
 S_{h,N,L}=\sum_{j=1}^{L}\frac{g_{N+h+j}-g_{N+j}}{2^{\,j}},\qquad
 R_{h,N,L}(M)=\sum_{j>L}\frac{M(N+h+j)+M(N+j)}{2^{\,j}} .
\]
If for every $h\ge1$ and every $N_0$ there are $N\ge N_0$ and $L\ge1$ with
$\operatorname{dist}(S_{h,N,L},\Z)>R_{h,N,L}(M)$, then the universal shift
escape of Problem~\ref{long251:prob:escape} holds.
\end{proposition}

\begin{proof}
The part of $\sum_{j\ge1}(g_{N+h+j}-g_{N+j})2^{-j}$ omitted from $S_{h,N,L}$
has absolute value at most $R_{h,N,L}(M)$, so under the displayed inequality
the full sum lies at positive distance from every integer.
\end{proof}

The truncation is an exact modular small-arc problem.  Writing the integral
dyadic block
\[
 D_{h,N,L}=\sum_{j=1}^{L}2^{\,L-j}(g_{N+h+j}-g_{N+j}),
 \qquad S_{h,N,L}=\frac{D_{h,N,L}}{2^{L}},
\]
one has
\[
 \operatorname{dist}(S_{h,N,L},\Z)
 =2^{-L}\min\bigl\{D_{h,N,L}\bmod2^{L},\;
 2^{L}-(D_{h,N,L}\bmod2^{L})\bigr\},
\]
where $D_{h,N,L}\bmod2^{L}$ is the least nonnegative residue, including when
$D_{h,N,L}<0$.  The finite criterion therefore asks the residue of
$D_{h,N,L}$ to avoid the two arcs of radius $2^{L}R_{h,N,L}(M)$ around $0$
modulo $2^{L}$.  A one-block certificate would be a prime-gap theorem
producing such an avoided arc on a logarithmic block.  This rewriting is
paper-level and carries no Lean declaration.

\begin{proposition}[completeness of finite separation]\label{long251:res:complete-truncation}
Suppose $D\in\R$, $S_L\in\R$ and $R_L\ge0$ satisfy
$|D-S_L|\le R_L$ and $R_L\to0$.  Then
\[
 D\notin\Z\quad\Longleftrightarrow\quad
 \text{there exists }L\text{ with }\operatorname{dist}(S_L,\Z)>R_L.
\]
\end{proposition}
\begin{proof}
The reverse implication follows from the triangle inequality.
For the forward implication put $\delta=\operatorname{dist}(D,\Z)>0$
and choose $R_L<\delta/2$.  The distance to $\Z$ is $1$-Lipschitz, so
$\operatorname{dist}(S_L,\Z)\ge\delta-R_L>R_L$.
\end{proof}
Thus existence of an arbitrarily deep finite certificate is an exact
reformulation of nonintegrality, not a new prime-gap estimate.  A useful
analytic producer must control the separation as well as the tail.

\subsection{Divisor-hitting shift escape}\label{long251:sec:xr-divisorhit}

\begin{problem}[divisor-hitting shift escape]\label{long251:xr:divisor-hit}\compatlabel{prob:divisor-hit}
For every $r\ge1$, does some positive multiple of $r$ escape cofinally, in the
sense of \eqref{long251:eq:shift-escape} at shift length $mr$ for some $m\ge1$?
\end{problem}

This is an exact reformulation, not a weaker mathematical target.
If $T_0$ is irrational, every positive shift is nonintegral by the block
identity, so divisor-hitting escape follows with $m=1$.  Conversely,
rationality gives a shift $h\ge1$ that is integral at every sufficiently
late index.  For each positive integer $m$, the identity
\[
 T_{N+mh}-T_N=\sum_{j=0}^{m-1}(T_{N+(j+1)h}-T_{N+jh})
\]
makes every multiple $mh$ eventually integral.  Taking $r=h$ contradicts
divisor-hitting escape.  This ordinary deduction uses the checked
rationality classification and a finite telescope; rearranging the
quantifiers supplies no new estimate for the actual prime gaps.

\subsection{Two ordinary proofs about the free-pair producer}
\label{long251:sec:xr-compression}

Both arguments below are ordinary mathematical deductions.  The first uses
the prime number theorem, while the second consumes the free-pair lattice
proved above.  Neither deduction is claimed to be kernel-checked; the exact
tail recurrence and free-pair statements they use have their separate linked
Lean evidence.

\paragraph{State compression.}  Exchanging the order of summation gives
$\sum_{N\le X}T_N\le(p_{X+1}-p_0)+2T_X$.  The prime number theorem
\cite[Ch.~6]{mv2007} also gives
\begin{equation}\label{long251:eq:tail-small-relative-prime}
 \frac{T_X}{p_X}
 =\sum_{j\ge1}\frac{g_{X+j}}{2^jp_X}\longrightarrow0.
\end{equation}
Indeed, for each fixed $j$ the summand without $2^{-j}$ tends to zero because
$p_{X+j}/p_X\to1$.  The usual two-sided bounds $p_n\asymp n\log n$ dominate
$g_{X+j}/p_X$ by $K(1+j)^2$, independently of $X$; multiplication by $2^{-j}$
therefore permits dominated convergence.  Applying the first-moment bound at
$2X$ and using \eqref{long251:eq:tail-small-relative-prime}, there is an absolute
$K_0$ such that
\[
 \sum_{X<N\le2X}T_N\le K_0X\log X
\]
for all large $X$.  Thus, for each fixed $C>1$, Markov's inequality leaves at
least $(1-1/C)X$ indices in $(X,2X]$ with $T_N\le K_0C\log X$.

Under rationality, take $X$ beyond the point where the reduced denominator is
the fixed odd integer $d$.  On this set $T_N$ lies in $d^{-1}\Z_{\ge0}$ and
takes at most $dK_0C\log X+1$ values.  Pigeonhole therefore gives one value at
$\gg_{C,d}X/\log X$ indices.  This compression uses the growth of the actual
primes.  It is false for a general rational integer-digit orbit: the quadratic
countermodel of Proposition~\ref{long251:res:polynomialcountermodel} has the
strictly increasing tails $T_N=2(N+4)^2$.  Even for the primes, equal-tail
pairs have difference zero and hence do not supply the nonintegral difference
required by the free-pair criterion.  The naive form
$\sum_{N\le X}T_N\le p_{X+1}$ is false.

\paragraph{Redundancy of the gap condition.}  Under the free-pair lattice the
difference $T_M-T_N$ is an integer whenever the modulus divides the offset, and
an integer cannot lie in $(\tfrac12,1)$.  The half-window component of the
reduced event of Theorem~\ref{long251:res:signedwindow} is therefore already a
contradiction on its own, and the gap-mismatch condition adds nothing to it.
That changes which component a proof has to produce.

\subsection{A bounded companion to the recurring-values countermodel}
\label{long251:sec:xr-boundedpolignac}

Theorem~\ref{long251:res:polignacfail} carries a logarithmic growth profile so that its
cumulative sequence matches the primes.  The mechanism is visible in a much
smaller example, which is recorded here because it is checkable by hand.

\begin{proposition}[bounded recurring-values countermodel]\label{long251:xr:boundedpolignac}
Put $U_0=4$ and, for $n\ge1$, $U_n=6$ when $n=k!$ for some $k\ge3$ and $U_n=4$
otherwise, and set $a_n=2U_{n-1}-U_n$ for $n\ge1$.  Then $a_n\in\{2,4,8\}$,
the value $2$ occurs at every index $k!$ and the value $4$ at every index
$2\,k!$, so both recur infinitely often at indices divisible by any fixed
$t\ge1$.  The series $\sum_{n\ge1}a_n2^{-n}$ equals $4$ and every tail
$\sum_{j\ge1}a_{N+j}2^{-j}$ equals the integer $U_N$.
\end{proposition}

\begin{proof}
For $k\ge3$ the index $k!$ is even and at least $6$, and $k!-1$ is odd, so the
predecessor of a spike is an ordinary index; hence $a_{k!}=8-6=2$,
$a_{k!+1}=12-4=8$, and $a_n=8-4=4$ elsewhere.  For $k\ge2$ one has
$k!<2\,k!<(k+1)!$, so $2\,k!$ is not a factorial, and $2\,k!-1$ is odd so it is
not a factorial either; therefore $a_{2k!}=4$.  Since $t\mid k!$ for every
large $k$, both values recur in the residue class $0$ modulo $t$.  The identity
$a_n2^{-n}=U_{n-1}2^{-(n-1)}-U_n2^{-n}$ telescopes to
$\sum_{j=1}^{m}a_{N+j}2^{-j}=U_N-U_{N+m}2^{-m}$, and $U$ is bounded, so the
endpoint vanishes.
\end{proof}

The cumulative sequence of this word grows linearly, which is the one property
that separates it from the primes and the reason the note carries the
logarithmic version instead.  Both refute the same generic implication.

\subsection{The measured densities in full}\label{long251:sec:xr-densities}

The note reports the range of the adjacent-mismatch density.  The per-offset
figures below come from the same scan over the $6\,841\,648$ primes under
$1.2\times10^{8}$, with double-precision tails.

\begin{center}
\small
\begin{tabular}{@{}rrrrrr@{}}
\toprule
$h$ & events & density & $\delta=+2$ & $\delta=-2$ & last event at prime\\
\midrule
1 & 56\,427 & 0.008248 & 28\,022 & 28\,405 & 119\,995\,753\\
2 & 31\,979 & 0.004674 & 15\,742 & 16\,237 & 119\,993\,807\\
3 & 30\,233 & 0.004419 & 15\,093 & 15\,140 & 119\,998\,321\\
4 & 29\,262 & 0.004277 & 14\,606 & 14\,656 & 119\,993\,473\\
\bottomrule
\end{tabular}
\end{center}

Every offset $h\le16$ has events, with minimum density $0.00418$ and maximum
$0.008248$; the two signs are near balanced at every offset.  For $h=1$ the
eight band densities, and the same figures after rescaling by the band mean of
$\log p$, run
\[
\begin{array}{l}
 0.011285,\;0.008951,\;0.008303,\;0.007936,\;0.007651,\;0.007507,\;0.007253,\;0.007094;\\[2pt]
 0.17671,\;0.15058,\;0.14420,\;0.14067,\;0.13766,\;0.13667,\;0.13333,\;0.13148 .
\end{array}
\]
The rescaled drift is $0.744$ at $h=1$ and lies between $0.744$ and $0.8121$
across all offsets.  Theorem~\ref{long251:res:sparse} proves that the limiting density
is zero, so the observed decline is the expected behaviour rather than a
failure of the event.

\subsection{Corrections to earlier records}\label{long251:sec:xr-corrections}

\paragraph{Two different improvements to the exclusion record.}
An earlier internal record gave a denominator exclusion at $10^{602}$.
The kernel-decided floor $2^{589}>10^{177}$ of
Theorem~\ref{long251:res:denominatorfloor} is numerically smaller but has a
different verification status.  It does not supersede $10^{602}$ in
numerical strength.  The continued-fraction exclusion
$2^{39997}>10^{12040}$ of Theorem~\ref{long251:res:cfexclusion} is the larger
numerical bound.  Its receipt's digit count $12041$ must not be substituted
for the exponent $12040$.  The independent integer replay supplied with
this revision reproduces the stated larger exclusion by a separate
interval and Farey-neighbour calculation; it is not a new Lean replay.

\paragraph{The named missing input.}  An earlier record named
Hardy-Littlewood $k$-tuple correlation of consecutive gaps at a fixed offset as
the missing input.  With the offset freed by Theorem~\ref{long251:res:freepair} the
route needs, for each modulus $t$, cofinally many congruent pairs with
certified nonintegral tail difference.  A further record proposed that two even
values differing by $2$, each occurring infinitely often inside a common index
residue class, would supply the two-condition form.  No proof of that
implication was ever produced, and Theorem~\ref{long251:res:polignacfail} shows that no
proof exists at the level of integer-digit recurrences.  That entry is
withdrawn.

\paragraph{The formal boundary.}  An earlier reading of the short note placed
the real-tail bridge and the real form of the adjacent-pair consumer on the
paper side.  Both are kernel-checked, at
\href{https://github.com/wcook04/plectis-erdos/blob/3d6d938d696fed0fb71dd55115a18a73738ff223/lean/ErdosProblems/Erdos251/RealPrimeGapTail.lean\#L48}{line~48}
and
\href{https://github.com/wcook04/plectis-erdos/blob/3d6d938d696fed0fb71dd55115a18a73738ff223/lean/ErdosProblems/Erdos251/RealPrimeGapTail.lean\#L99}{line~99}
of the real prime-gap tail module.  A second reading placed the kernel-decided
denominator floor and the free-pair criterion outside the source revision the
note pins.  Both modules are present at that revision, and both are linked from
the note.
% ---- part family_catalogue ----
\clearpage
\section{Complete affine-lattice statement from the short note}
\label{long251:app:affine-detail}
The following expands the compact circularity statement formerly printed
in the body of this record.  In particular the fixed-lattice predicate
and its strict growth hypothesis are stated, not left implicit.

Here $D_N=\sigma_h(N)$ and
$\delta_N=g_{N+h+1}-g_{N+1}$, with $h$ fixed.  For a longer block, set
\[
 B_{h,N,r}=\sum_{i=0}^{r-1}2^{r-1-i}\delta_{N+i};
 \qquad D_{N+r}=2^rD_N-B_{h,N,r}.
\]
Call the data-dependent affine condition
\[
 \mathcal A_{N,r}:\qquad
 D_{N+r}\in -B_{h,N,r}+2^{r+1}\Z.
\]

\begin{theorem}[affine and fixed-lattice circularity]\label{long251:res:affinecollapse}\compatlabel{res:affinecollapse}
For every rational dyadic tail recurrence and all $h,N,r\ge0$,
\begin{equation}
 \mathcal A_{N,r}\quad\Longleftrightarrow\quad D_N\in2\Z.
\label{eq:affinecollapse}\compatlabel{long251:eq:affinecollapse}
\end{equation}
Consequently, if every $\delta_N$ is even, then
\begin{equation}
 \bigl(\forall N_0\ \exists N,r:\ N_0<N\text{ and }\neg\mathcal A_{N,r}\bigr)
 \quad\Longleftrightarrow\quad
 D_N\notin\Z\text{ for arbitrarily large }N.
\label{eq:affinecofinal}\compatlabel{long251:eq:affinecofinal}
\end{equation}

There is a second equivalence.  Let $b:\N\to\Q$ satisfy
$|D_N|\le b(N)$ for every $N$, and suppose that for every $N$ and every
positive integer $q$ there is an $r$ with
\begin{equation}
 2b(N+r)q<2^r.
\label{eq:dyadicscale}\compatlabel{long251:eq:dyadicscale}
\end{equation}
Then
\begin{equation}
 \begin{split}
 &\forall N_0\ \exists N,r:\ N_0<N\text{ and }
   \forall z\in\Z,\quad
   b(N+r)<|B_{h,N,r}-2^rz|\\
 &\hspace{35mm}\Longleftrightarrow\quad
 D_N\notin\Z\text{ for arbitrarily large }N.
 \end{split}
\label{eq:fixedcollapse}\compatlabel{long251:eq:fixedcollapse}
\end{equation}
\end{theorem}

\begin{proof}
Substituting $D_{N+r}=2^rD_N-B_{h,N,r}$ into $\mathcal A_{N,r}$ cancels the
observed block from both sides and leaves $D_N=2z$.  This proves
\eqref{long251:eq:affinecollapse}.  After one recurrence step, evenness of
$\delta_N$ identifies even integrality at $N+1$ with ordinary integrality at
$N$.  The forward implication in~\eqref{long251:eq:affinecofinal} follows, while the
reverse implication chooses a nonintegral $D_N$ and the legal depth $r=0$.

For~\eqref{long251:eq:fixedcollapse}, eventual integrality and the block identity give
some $z\in\Z$ with
\[
 B_{h,N,r}-2^rz=-D_{N+r},
\]
so the displayed strict separation contradicts $|D_{N+r}|\le b(N+r)$.
Conversely, if $D_N$ is nonintegral and has reduced denominator $q$, then its
distance from every integer is at least $1/q$.  Choose $r$ from
\eqref{long251:eq:dyadicscale}, scale this separation by $2^r$, and use the block
identity together with $|D_{N+r}|\le b(N+r)$.  The triangle inequality gives
$|B_{h,N,r}-2^rz|>b(N+r)$ for every integer $z$.
\end{proof}

\noindent The three displayed equivalences are assembled in one declaration,
\href{https://github.com/wcook04/plectis-erdos/blob/3d6d938d696fed0fb71dd55115a18a73738ff223/lean/ErdosProblems/Erdos251/PaperCoreR7.lean\#L241}{affine circularity bundle}.

The first equivalence shows that the apparent affine hierarchy contains no
depth-dependent information: it is the pullback of even integrality through
the recurrence identity.  The fixed lattice removes that data dependence, but
under the stated growth condition rational denominator separation already
forces every required escape.  Hence neither cofinal condition is an
independent source of information about consecutive primes.

\section{Mathematical and evidence map}\label{sec:erdos-251-complete-family-map}
\paragraph{Perturbations.}
The bounded construction and the sparse congruence-preserving theorem are
existence results for altered integer words.  Their conclusion is not a
prime-producing construction.  The ordinary sparse proof and the formal
existence proof have different witnesses.
\paragraph{Exact criteria.}
Summation by parts identifies the actual tails.  Integral-shift, free-pair
and lcm-diagonal criteria then classify rationality at the recurrence level.
Their equivalence does not supply the missing prime-specific witnesses.
\paragraph{Obstructions.}
Polynomial growth, bounded residues, fixed-block nonconcentration and
recurring gap values do not by themselves force the needed tail behaviour.
The numbered constructions above give the precise counterexamples.
\paragraph{Certificates.}
Finite window comparisons and denominator exclusions are exact finite
statements once the proved analytic remainder bounds are supplied.  A
large rational-denominator floor is not an irrationality theorem.
\paragraph{Formal evidence.}
The revision's declaration ledger records the file, line, immutable
revision and supplied build status of each linked declaration.  The older
administrative family catalogue is retained verbatim with that ledger;
counts of registry rows or project-wide families do not strengthen any
mathematical statement.

% ---- part back ----
\begin{thebibliography}{99}
\small\raggedright\setlength{\itemsep}{0pt}
\bibitem{erdos1958} Paul Erd\H{o}s, \href{https://users.renyi.hu/~p_erdos/1958-19.pdf}{\emph{Sur certaines s\'eries \`a valeur irrationnelle}}.
  L'Enseignement Mathématique \textbf{4} (1958), 93--100, doi:\href{https://doi.org/10.5169/seals-34629}{10.5169/seals-34629}.
\bibitem{erdosgraham1980} Paul Erd\H{o}s and Ronald L. Graham, \href{https://mathweb.ucsd.edu/~ronspubs/80_11_number_theory.pdf}{\emph{Old and New Problems and Results in Combinatorial Number Theory}}.
  Monographies de L'Enseignement Mathématique, L'Enseignement Mathématique, 1980.
\bibitem{erdospomerance1978} Paul Erd\H{o}s and Carl Pomerance, \href{https://doi.org/10.1007/BF01818569}{\emph{On the largest prime factors of $n$ and $n+1$}}.
  Aequationes Mathematicae \textbf{17} (1978), 311--321, doi:\href{https://doi.org/10.1007/BF01818569}{10.1007/BF01818569}.
\bibitem{erdos1988} Paul Erd\H{o}s, \href{https://doi.org/10.1017/CBO9780511897184.009}{\emph{On the irrationality of certain series: problems and results}}.
  New Advances in Transcendence Theory, Cambridge University Press, 1988, pp.~102--109, doi:\href{https://doi.org/10.1017/CBO9780511897184.009}{10.1017/CBO9780511897184.009}.
\bibitem{erdos1932} Paul Erd\H{o}s, \href{https://users.renyi.hu/~p_erdos/1932-01.pdf}{\emph{Beweis eines Satzes von Tschebyschef}}.
  Acta Litterarum ac Scientiarum Szeged \textbf{5} (1932), 194--198.
\bibitem{erdosstraus1963} Paul Erd\H{o}s and Ernst G. Straus, \emph{On the irrationality of certain Ahmes series}.
  Journal of the Indian Mathematical Society (N.S.) \textbf{27} (1964), 129--133. Legacy citation key retained; supplied bibliography gives publication year 1964. MR0175848.
\bibitem{erdosstraus1974} Paul Erd\H{o}s and Ernst G. Straus, \href{https://doi.org/10.2140/pjm.1974.55.85}{\emph{On the irrationality of certain series}}.
  Pacific Journal of Mathematics \textbf{55} (1974), no.~1, 85--92, doi:\href{https://doi.org/10.2140/pjm.1974.55.85}{10.2140/pjm.1974.55.85}.
\bibitem{fgkmt2018} Kevin Ford, Ben Green, Sergei Konyagin, James Maynard and Terence Tao, \href{https://arxiv.org/abs/1412.5029v3}{\emph{Long gaps between primes}}.
  Journal of the American Mathematical Society \textbf{31} (2018), 65--105, doi:\href{https://doi.org/10.1090/jams/876}{10.1090/jams/876}; arXiv:\href{https://arxiv.org/abs/1412.5029v3}{1412.5029v3}.
\bibitem{fridy1966} John A. Fridy, \href{https://www.fq.math.ca/Scanned/4-3/fridy.pdf}{\emph{Generalized bases for the real numbers}}.
  The Fibonacci Quarterly \textbf{4} (1966), no.~3, 193--201.
\bibitem{khinchin1964} A. Ya. Khinchin, \emph{Continued Fractions}.
  University of Chicago Press, 1964.
\bibitem{kovactao2024} Vjekoslav Kova\v{c} and Terence Tao, \href{https://arxiv.org/abs/2406.17593v4}{\emph{On several irrationality problems for Ahmes series}}.
  Acta Mathematica Hungarica \textbf{175} (2025), 572--608, doi:\href{https://doi.org/10.1007/s10474-025-01528-0}{10.1007/s10474-025-01528-0}; arXiv:\href{https://arxiv.org/abs/2406.17593v4}{2406.17593v4}. Statement and section numbers refer to arXiv version 4; legacy citation key retained.
\bibitem{kuperberg2023} Vivian Kuperberg, \href{https://arxiv.org/abs/2210.09775v2}{\emph{Sums of singular series with large sets and the tail of the distribution of primes}}.
  The Quarterly Journal of Mathematics \textbf{74} (2023), no.~4, 1457--1479, doi:\href{https://doi.org/10.1093/qmath/haad030}{10.1093/qmath/haad030}; arXiv:\href{https://arxiv.org/abs/2210.09775v2}{2210.09775v2}. Statement numbers refer to arXiv version 2, 15 June 2023.
\bibitem{land2026} Johan Land, \href{https://github.com/beetree/math_erdos_251}{\emph{A conditional proof of the irrationality of $\sum_{n\ge1}p_n2^{-n}$ under a uniform Hardy--Littlewood prime-tuples conjecture}}.
  Research draft, 5 September 2026, 2026. The author supplies formalisation material. No independent rebuild is claimed in this revision.
\bibitem{lean4} Leonardo de Moura and Sebastian Ullrich, \href{https://doi.org/10.1007/978-3-030-79876-5_37}{\emph{The Lean 4 theorem prover and programming language}}.
  Automated Deduction -- CADE 28, Lecture Notes in Computer Science, Springer, 2021, pp.~625--635, doi:\href{https://doi.org/10.1007/978-3-030-79876-5_37}{10.1007/978-3-030-79876-5\_37}.
\bibitem{mathlib} {The mathlib Community}, \href{https://doi.org/10.1145/3372885.3373824}{\emph{The Lean mathematical library}}.
  Proceedings of the 9th ACM SIGPLAN International Conference on Certified Programs and Proofs, ACM, 2020, pp.~367--381, doi:\href{https://doi.org/10.1145/3372885.3373824}{10.1145/3372885.3373824}. Describes a Lean 3-era snapshot; the repository lock identifies the library used by the present Lean 4 sources.
\bibitem{maynard2015} James Maynard, \href{https://doi.org/10.4007/annals.2015.181.1.7}{\emph{Small gaps between primes}}.
  Annals of Mathematics \textbf{181} (2015), 383--413, doi:\href{https://doi.org/10.4007/annals.2015.181.1.7}{10.4007/annals.2015.181.1.7}.
\bibitem{mv2007} Hugh L. Montgomery and Robert C. Vaughan, \href{https://doi.org/10.1017/CBO9780511618314}{\emph{Multiplicative Number Theory I: Classical Theory}}.
  Cambridge Studies in Advanced Mathematics, Cambridge University Press, 2007, doi:\href{https://doi.org/10.1017/CBO9780511618314}{10.1017/CBO9780511618314}.
\bibitem{schlagepuchta2011} Jan-Christoph Schlage-Puchta, \href{https://arxiv.org/abs/1105.1451v1}{\emph{The irrationality of some number theoretical series}}.
  Acta Arithmetica \textbf{126} (2007), no.~4, 295--303, doi:\href{https://doi.org/10.4064/aa126-4-1}{10.4064/aa126-4-1}; arXiv:\href{https://arxiv.org/abs/1105.1451v1}{1105.1451v1}. Published in 2007; arXiv upload is from 2011. Lemma and preprint page locators refer to arXiv version 1; legacy citation key retained.
\bibitem{short2009} Ian Short, \href{https://arxiv.org/abs/0912.1997v1}{\emph{Ford Circles, Continued Fractions, and Rational Approximation}}.
  The American Mathematical Monthly \textbf{118} (2011), no.~2, 130--135, doi:\href{https://doi.org/10.4169/amer.math.monthly.118.02.130}{10.4169/amer.math.monthly.118.02.130}; arXiv:\href{https://arxiv.org/abs/0912.1997v1}{0912.1997v1}. Preprint title: Ford circles, continued fractions, and best approximation of the second kind. Preprint theorem locators refer to arXiv version 1; legacy citation key retained.
\bibitem{vandoornkovac2025} Wouter van Doorn and Vjekoslav Kova\v{c}, \href{https://arxiv.org/abs/2509.24971v3}{\emph{Lacunary sequences whose reciprocal sums represent all rational numbers in an interval}}.
  Acta Arithmetica \textbf{223} (2026), 275--295, doi:\href{https://doi.org/10.4064/aa251001-13-1}{10.4064/aa251001-13-1}; arXiv:\href{https://arxiv.org/abs/2509.24971v3}{2509.24971v3}. Statement numbers refer to arXiv version 3; legacy citation key retained.
\bibitem{zhang2014} Yitang Zhang, \href{https://doi.org/10.4007/annals.2014.179.3.7}{\emph{Bounded gaps between primes}}.
  Annals of Mathematics \textbf{179} (2014), 1121--1174, doi:\href{https://doi.org/10.4007/annals.2014.179.3.7}{10.4007/annals.2014.179.3.7}.
\bibitem{erdosproblems} Thomas F. Bloom, \href{https://www.erdosproblems.com/251}{\emph{Erd\H{o}s Problem \#251}}.
  2026. Catalogue snapshot inherited from the supplied manuscript, accessed 6 September 2026. Live page was inaccessible to this revision on 16 September 2026; no refreshed global-status claim is inferred.
\bibitem{erdosproblems251thread} {Erdős Problems contributors}, \href{https://www.erdosproblems.com/forum/thread/251}{\emph{Erd\H{o}s Problem \#251 discussion thread}}.
  2026. Inherited snapshot accessed 6 September 2026: Tao comment of 7 October 2025 and Land comments of 6 September 2026. Live thread not reverified.
\bibitem{kovac2026} {ChatGPT 5.4 Pro (orchestrated by Vjeko Kova\v{c})}, \href{https://web.math.pmf.unizg.hr/~vjekovac/files/Erdos_problem_251.pdf}{\emph{On the Erd\H{o}s problem \#251}}.
  Unpublished note, Department of Mathematics, University of Zagreb, 2026. Printed attribution retained from the supplied bibliography; inherited access date 6 September 2026.
\bibitem{formalconjectures251} {The Formal Conjectures Authors}, \href{https://github.com/google-deepmind/formal-conjectures/blob/f776d2f2039351b00737ffcafb9d7d7666e1d9af/FormalConjectures/ErdosProblems/251.lean}{\emph{FormalConjectures.ErdosProblems.251}}.
  2025. Pinned statement source, not a proof of Problem 251; inherited access date 28 July 2026.
\bibitem{ringer2026} Stefan Ringer, \href{https://github.com/StefanRinger/erdos-251}{\emph{Local gap statistics, telescoping, and normality: a local-pattern approach to Erd\H{o}s problem 251}}.
  Preprint dated 11 September 2026, 2026. Mutable main-branch TeX consulted 16 September 2026. The author supplies formalisation material; this revision does not independently rebuild it.
\bibitem{crmarickovac2025} Ton\'{c}i Crmari\'{c} and Vjekoslav Kova\v{c}, \href{https://arxiv.org/abs/2504.18712v1}{\emph{On the irrationality of certain super-polynomially decaying series}}.
  Colloquium Mathematicum \textbf{179} (2025), 55--68, doi:\href{https://doi.org/10.4064/cm9628-5-2025}{10.4064/cm9628-5-2025}; arXiv:\href{https://arxiv.org/abs/2504.18712v1}{2504.18712v1}. Lemma 4 is cited using arXiv version 1.
\bibitem{adm2022} Boris Adamczewski, Michael Drmota and Clemens M\"ullner, \href{https://arxiv.org/abs/2009.14773v2}{\emph{(Logarithmic) densities for automatic sequences along primes and squares}}.
  Transactions of the American Mathematical Society \textbf{375} (2022), no.~1, 455--499, doi:\href{https://doi.org/10.1090/tran/8476}{10.1090/tran/8476}; arXiv:\href{https://arxiv.org/abs/2009.14773v2}{2009.14773v2}. Theorems 1.2 and 1.4 refer to arXiv version 2, 13 April 2021; journal publication is 2022.
\bibitem{vandoorn2025} Wouter van Doorn, \href{https://arxiv.org/abs/2502.02200v2}{\emph{Partitions with prescribed sum of reciprocals: asymptotic bounds}}.
  2025; arXiv:\href{https://arxiv.org/abs/2502.02200v2}{2502.02200v2}. Version 2, 23 July 2025. No journal publication was confirmed in this pass.
\bibitem{glabprus2025} Szymon G{\l}\k{a}b and Franciszek Prus-Wi\'sniowski, \href{https://arxiv.org/abs/2512.17285v1}{\emph{Achievement sets -- current results and open problems}}.
  Real Analysis Exchange (2026), doi:\href{https://doi.org/10.14321/realanalexch.1766383782}{10.14321/realanalexch.1766383782}; arXiv:\href{https://arxiv.org/abs/2512.17285v1}{2512.17285v1}. Advance publication, first available in Project Euclid 8 June 2026. Consulted text remains arXiv:2512.17285v1 (19 December 2025); no final volume or page range was confirmed. Legacy citation key retained.
\bibitem{prusptak2024} Franciszek Prus-Wi\'sniowski and Jolanta Ptak, \href{https://arxiv.org/abs/2412.08768v1}{\emph{Achievable Cantorvals almost without reversed Kakeya conditions}}.
  2024; arXiv:\href{https://arxiv.org/abs/2412.08768v1}{2412.08768v1}. Version 1 submitted 11 December 2024. The sparse indices satisfy the overlap inequality, not its strict term-dominating reverse.
\bibitem{bartoszewicz2014} Artur Bartoszewicz, Ma{\l}gorzata Filipczak and Emilia Szymonik, \href{https://arxiv.org/abs/1304.4218v2}{\emph{Multigeometric sequences and Cantorvals}}.
  Central European Journal of Mathematics \textbf{12} (2014), no.~7, 1000--1007, doi:\href{https://doi.org/10.2478/s11533-013-0396-4}{10.2478/s11533-013-0396-4}; arXiv:\href{https://arxiv.org/abs/1304.4218v2}{1304.4218v2}.
\bibitem{kuperbergweighted2025} Vivian Kuperberg, \href{https://arxiv.org/abs/2301.06095v1}{\emph{Sums of singular series along arithmetic progressions and with smooth weights}}.
  International Journal of Number Theory \textbf{21} (2025), no.~1, 53--74, doi:\href{https://doi.org/10.1142/S1793042125500046}{10.1142/S1793042125500046}; arXiv:\href{https://arxiv.org/abs/2301.06095v1}{2301.06095v1}. Preprint uploaded 15 January 2023; journal publication is 2025.
\bibitem{jha2026} Abhishek Jha, \href{https://arxiv.org/abs/2605.23014v2}{\emph{The Poisson Tail Conjecture for primes in short intervals}}.
  2026; arXiv:\href{https://arxiv.org/abs/2605.23014v2}{2605.23014v2}. Substantially revised version 2, 12 September 2026, 31 pages. Conditional statements must retain their strong Hardy--Littlewood hypotheses.
\bibitem{polymath2014} {D. H. J. Polymath}, \href{https://arxiv.org/abs/1407.4897v4}{\emph{Variants of the Selberg sieve, and bounded intervals containing many primes}}.
  Research in the Mathematical Sciences \textbf{1} (2014), article 12, doi:\href{https://doi.org/10.1186/s40687-014-0012-7}{10.1186/s40687-014-0012-7}; arXiv:\href{https://arxiv.org/abs/1407.4897v4}{1407.4897v4}. Theorem 1.4(i) refers to arXiv version 4 (22 December 2014); the journal numbers it Theorem 4(i). An erratum is recorded at doi:10.1186/s40687-015-0033-x.
\bibitem{nitecki2015} Zbigniew Nitecki, \href{https://arxiv.org/abs/1106.3779v2}{\emph{Cantorvals and Subsum Sets of Null Sequences}}.
  The American Mathematical Monthly \textbf{122} (2015), no.~9, 862--870, doi:\href{https://doi.org/10.4169/amer.math.monthly.122.9.862}{10.4169/amer.math.monthly.122.9.862}; arXiv:\href{https://arxiv.org/abs/1106.3779v2}{1106.3779v2}. Consulted preprint: Subsum Sets: Intervals, Cantor Sets, and Cantorvals, version 2 (8 July 2013). Theorem 14 is attributed there to Guthrie--Nymann; its locator is not journal pagination.
\bibitem{miskaprusptak2023} Piotr Miska, Franciszek Prus-Wi\'sniowski and Jolanta Ptak, \href{https://ruj.uj.edu.pl/server/api/core/bitstreams/d6630f7b-e6ee-4de8-8a1b-81c7b4c59d2e/content}{\emph{More on Kakeya Conditions for Achievement Sets}}.
  Results in Mathematics \textbf{78} (2023), article 113, doi:\href{https://doi.org/10.1007/s00025-023-01890-x}{10.1007/s00025-023-01890-x}. Repairs an estimate in the 2021 proof, preserving its uniqueness conclusion, and gives a simpler proof of a weaker theorem without that conclusion.
\bibitem{marchwickimiska2021} Jacek Marchwicki and Piotr Miska, \href{https://link.springer.com/article/10.1007/s00025-021-01479-2}{\emph{On Kakeya Conditions for Achievement Sets}}.
  Results in Mathematics \textbf{76} (2021), article 181, doi:\href{https://doi.org/10.1007/s00025-021-01479-2}{10.1007/s00025-021-01479-2}. Theorem 2.1 is to be read with the proof repair in Miska--Prus-Wi\'sniowski--Ptak (2023).
\bibitem{nowakowski2025} Piotr Nowakowski, \href{https://arxiv.org/abs/2512.17761v1}{\emph{On a new condition implying that an achievement set is a Cantorval and its applications}}.
  2025; arXiv:\href{https://arxiv.org/abs/2512.17761v1}{2512.17761v1}. Version 1, 19 December 2025. Theorem 3.1 requires the Star Procedure of Definition 2 never to break; no application to the present factorial weights is asserted.
\end{thebibliography}

\end{document}
